% concatenated from fesddr.tex, definitionsSym.tex, definitionsEnv.tex
\documentclass[10pt,a4paper]{article}
%\usepackage[utf8]{inputenc}
%\usepackage[T1]{fontenc}
\usepackage{amsmath,amssymb,amsfonts,amsthm}
%\usepackage{mathabx,MnSymbol}
%\usepackage{graphicx}
\usepackage{float,graphicx,color}
\usepackage[all,cmtip]{xy}
\usepackage{tikz}
\usetikzlibrary{matrix}

%-------------definitions--------------------------------


%-----------typo---------------------
%\newcommand{\hedera}{\Pisymbol{pzd}{166}} %or 167
%\newcommand{\pilcrow}{\Pisymbol{pzd}{161}}
%\newcommand{\malta}{\Pisymbol{pzd}{32}}

\newcommand{\hedera}{\ding{166}}
\newcommand{\pilcrow}{\ding{161}}
\newcommand{\malta}{\ding{32}}
\newcommand{\dangersym}{$\red{\triangle}$}





%-----------------mathcharacters------


\newcommand{\bfr}{\mathbf{r}}

\newcommand{\rmd}{\mathrm{d}}
\newcommand{\rmr}{\mathrm{r}}

\newcommand{\rmA}{\mathrm{A}}
\newcommand{\rmB}{\mathrm{B}}
\newcommand{\rmC}{\mathrm{C}}
\newcommand{\rmD}{\mathrm{D}}
\newcommand{\rmH}{\mathrm{H}}
\newcommand{\rmI}{\mathrm{I}}
\newcommand{\rmL}{\mathrm{L}}
\newcommand{\rmN}{\mathrm{N}}
\newcommand{\rmO}{\mathrm{O}}
\newcommand{\rmP}{\mathrm{P}}
\newcommand{\rmR}{\mathrm{R}}
\newcommand{\rmS}{\mathrm{S}}
\newcommand{\rmT}{\mathrm{T}}
\newcommand{\rmU}{\mathrm{U}}
\newcommand{\rmW}{\mathrm{W}}

\newcommand{\sfd}{\mathsf{d}}

\newcommand{\bbd}{\mathbb{d}}

\newcommand{\bbA}{\mathbb{A}}
\newcommand{\bbB}{\mathbb{B}}
\newcommand{\bbC}{\mathbb{C}}
\newcommand{\bbD}{\mathbb{D}}
\newcommand{\bbE}{\mathbb{E}}
\newcommand{\bbF}{\mathbb{F}}
\newcommand{\bbG}{\mathbb{G}}
\newcommand{\bbH}{\mathbb{H}}
\newcommand{\bbI}{\mathbb{I}}
\newcommand{\bbJ}{\mathbb{J}}
\newcommand{\bbK}{\mathbb{K}}
\newcommand{\bbL}{\mathbb{L}}
\newcommand{\bbM}{\mathbb{M}}
\newcommand{\bbN}{\mathbb{N}}
\newcommand{\bbP}{\mathbb{P}}
\newcommand{\bbQ}{\mathbb{Q}}
\newcommand{\bbR}{\mathbb{R}}
\newcommand{\bbS}{\mathbb{S}}
\newcommand{\bbT}{\mathbb{T}}
\newcommand{\bbU}{\mathbb{U}}
\newcommand{\bbV}{\mathbb{V}}
\newcommand{\bbW}{\mathbb{W}}
\newcommand{\bbZ}{\mathbb{Z}}

\newcommand{\frd}{\mathfrak{d}}

\newcommand{\frg}{\mathfrak{g}}
\newcommand{\frk}{\mathfrak{k}}
\newcommand{\frl}{\mathfrak{l}}
\newcommand{\fro}{\mathfrak{o}}
\newcommand{\frp}{\mathfrak{p}}

\newcommand{\frs}{\mathfrak{s}}
\newcommand{\fru}{\mathfrak{u}}


\newcommand{\frA}{\mathfrak{A}}
\newcommand{\frE}{\mathfrak{E}}
\newcommand{\frH}{\mathfrak{H}}
\newcommand{\frI}{\mathfrak{I}}
\newcommand{\frK}{\mathfrak{K}}
\newcommand{\frL}{\mathfrak{L}}
\newcommand{\frP}{\mathfrak{P}}
\newcommand{\frS}{\mathfrak{S}}
\newcommand{\frU}{\mathfrak{U}}
\newcommand{\frV}{\mathfrak{V}}

\newcommand{\scrC}{\mathscr{C}}
\newcommand{\scrD}{\mathscr{D}}
\newcommand{\scrT}{\mathscr{T}}

\newcommand{\calA}{\mathcal{A}}
\newcommand{\calB}{\mathcal{B}}
\newcommand{\calC}{\mathcal{C}}
\newcommand{\calD}{\mathcal{D}}
\newcommand{\calE}{\mathcal{E}}
\newcommand{\calF}{\mathcal{F}}
\newcommand{\calG}{\mathcal{G}}
\newcommand{\calH}{\mathcal{H}}
\newcommand{\calI}{\mathcal{I}}
\newcommand{\calJ}{\mathcal{J}}
\newcommand{\calK}{\mathcal{K}}
\newcommand{\calL}{\mathcal{L}}
\newcommand{\calM}{\mathcal{M}}
\newcommand{\calN}{\mathcal{N}}
\newcommand{\calO}{\mathcal{O}}
\newcommand{\calP}{\mathcal{P}}
\newcommand{\calQ}{\mathcal{Q}}
\newcommand{\calR}{\mathcal{R}}
\newcommand{\calS}{\mathcal{S}}
\newcommand{\calT}{\mathcal{T}}
\newcommand{\calU}{\mathcal{U}}
\newcommand{\calV}{\mathcal{V}}
\newcommand{\calW}{\mathcal{W}}
\newcommand{\calX}{\mathcal{X}}
\newcommand{\calY}{\mathcal{Y}}
\newcommand{\calZ}{\mathcal{Z}}


\newcommand{\calg}{\mathcal{g}}

\newcommand{\half}{\frac{1}{2}}

\newcommand{\bs}{{\scriptscriptstyle{\bullet}}}


%-----------math operators--------------

\newcommand{\one}{\mathbf{1}}
\newcommand{\zero}{\mathbf{0}}


\newcommand{\transp}{\mathrm{T}}   %{\textsc{t}}
\newcommand{\scT}{\transp}

\newcommand{\herm}{\textsc{h}}
\newcommand{\comp}{\textsc{c}}
\newcommand{\rota}{\textsc{r}}
\newcommand{\conj}{\dagger}


\newcommand{\id}{\mathrm{id}}

\DeclareMathOperator{\pb}{\mathrm{pb}}
%\newcommand{\tr}{\mathrm{tr}}

\DeclareMathOperator{\Div}{div}
\renewcommand{\div}{\Div}
\DeclareMathOperator{\rot}{rot}
\DeclareMathOperator{\grad}{grad}
\DeclareMathOperator{\curl}{curl}
\DeclareMathOperator{\defo}{def}
\DeclareMathOperator{\airy}{airy}
\DeclareMathOperator{\osc}{osc}

\newcommand{\tensorsym}{\otimes_{\mathsf{S}}}
%\newcommand{\gradsym}{\grad_{\mathsf{S}}}
\newcommand{\curltcurl}{\curl \transp\, \curl}
\DeclareMathOperator{\hess}{hess}
\DeclareMathOperator{\twist}{twist}
\DeclareMathOperator{\sven}{sven}
\DeclareMathOperator{\belt}{belt}
\DeclareMathOperator{\cott}{cott}

\newcommand{\Rie}{\mathrm{Rm}}
\newcommand{\Ric}{\mathrm{Rc}}


\DeclareMathOperator{\Skew}{skew}
\DeclareMathOperator{\op}{op}
\DeclareMathOperator{\PV}{PV}
\DeclareMathOperator{\Jac}{Jac}

\DeclareMathOperator{\vect}{vect}
\DeclareMathOperator{\myspan}{span}
\newcommand{\lspan}{\mathrm{span}}


%\DeclareMathOperator{\ker}{ker}
\DeclareMathOperator{\ran}{ran}
\DeclareMathOperator{\img}{img}
\DeclareMathOperator{\coker}{coker}
\DeclareMathOperator{\dom}{dom}
\DeclareMathOperator{\codom}{codom}

\DeclareMathOperator{\re}{\mathfrak{Re}}
\DeclareMathOperator{\im}{\mathfrak{Im}}

\DeclareMathOperator{\card}{{\scriptstyle \#}}
\DeclareMathOperator{\diam}{diam}
\DeclareMathOperator{\vol}{vol}
\DeclareMathOperator{\dist}{dist}
\DeclareMathOperator{\supp}{supp}
\DeclareMathOperator{\meas}{meas}
\DeclareMathOperator{\sign}{sign}

\DeclareMathOperator{\tr}{tr}
\DeclareMathOperator{\st}{st}

\DeclareMathOperator{\SL}{SL}
\DeclareMathOperator{\SU}{SU}
\DeclareMathOperator{\SO}{SO}
%\DeclareMathOperator{\sl}{\frs \frl}
\DeclareMathOperator{\su}{\frs \fru}
\DeclareMathOperator{\so}{\frs \fro}

\DeclareMathOperator{\gt}{\mathrm{GT}}
\DeclareMathOperator{\pt}{\mathrm{pt}}
\DeclareMathOperator{\hol}{\mathrm{hol}}


\DeclareMathOperator{\graf}{graf}

\DeclareMathOperator{\Hom}{Hom}
\DeclareMathOperator{\End}{End}

\DeclareMathOperator{\bigo}{\mathrm{O}}
\DeclareMathOperator{\smallo}{\mathrm{o}}


\DeclareMathOperator{\tracem}{\textsc{tr}}
\newcommand{\diag}{\textsc{di}}
\DeclareMathOperator{\symm}{\textsc{sym}}
\DeclareMathOperator{\asymm}{\textsc{asym}}
\DeclareMathOperator{\skewm}{\textsc{skewm}}
\DeclareMathOperator{\skewv}{\textsc{skewv}}
\DeclareMathOperator{\soper}{\textsc{s}}
\DeclareMathOperator{\dev}{\textsc{dev}}


\renewcommand{\Re}{\re}
\renewcommand{\Im}{\im}


\newcommand{\catGro}{\mathrm{Groups}}
\newcommand{\catRin}{\mathrm{Rings}}
\newcommand{\catSet}{\mathrm{Sets}}
\newcommand{\catTop}{\mathrm{Tops}}



\newcommand{\ts}{\textstyle}
\newcommand{\ds}{\displaystyle}

\newcommand{\Tan}{{\textsc{t}}}
\newcommand{\Nor}{{\textsc{n}}}

\newcommand{\rec}{\mathrm{rec}}
\newcommand{\loc}{{loc}}

\newcommand{\inc}{{inc}}
%\newcommand{\ext}{{ext}}
\newcommand{\tot}{{tot}}

%\newcommand{\comp}{\mathrm{comp}}




\newcommand{\Ltan}{\rmL_{\scriptscriptstyle \Tan}}
\newcommand{\Ln}{\rmL_{\nor}}
\newcommand{\LE}{\rmL_{\bbE}}
\newcommand{\La}{\rmL_{\mathrm a}}


\newcommand{\WF}{\mathfrak{W}}

\newcommand{\Lip}{\mathrm{Lip}}

\newcommand{\Dtan}{\scrD_{\scriptscriptstyle \Tan}}
\newcommand{\Dn}{\scrD_{\nor}}
\newcommand{\Dx}{\scrD_{\times}}

\newcommand{\Dslit}{\scrD^{\mathrm{slit}}}
\newcommand{\Dtanslit}{\Dtan^{\mathrm{slit}}}

\newcommand{\jump}[1]{[\![ #1 ]\!]}
\newcommand{\jumpn}[1]{\jump{#1}_{\nor}}
\newcommand{\jumpx}[1]{\jump{#1}_{\times}}

\newcommand{\pinor}{\pi_{\nor}}
\newcommand{\pitan}{\pi_{\scriptscriptstyle \Tan}}

\newcommand{\taunor}{\tau_{\nor}}
\newcommand{\tautan}{\tau_{\scriptscriptstyle \Tan}}

\newcommand{\sdot}{{\scriptscriptstyle \bullet}}

\newcommand{\gamman}{\gamma_{\nor}}
\newcommand{\gammatan}{\gamma_{\scriptscriptstyle \Tan}}
\newcommand{\gammax}{\gamma_{\times}}

\newcommand{\Hgrad}{\rmH_{\grad}}
\newcommand{\Hdiv}{\rmH_{\div}}
\newcommand{\Hrot}{\rmH_{\rot}}
\newcommand{\Hd}{\rmH_{\rmd}}
\newcommand{\Htan}{\rmH_{\scriptscriptstyle \Tan}}
\newcommand{\Hx}{\rmH_{\times}}
\newcommand{\Hn}{\rmH_{\nor}}
\newcommand{\Hsk}{\rmH_{\Skew}}
\newcommand{\Hcurl}{\rmH_{\curl}}
\newcommand{\Hop}{\rmH_{\op}}
\newcommand{\HE}{\rmH_{\bbE}}

\newcommand{\rmdsubA}{\rmd_{\! A}}
\newcommand{\mydot}{\! \dot{\phantom{A}} }

\newcommand{\LSF}{\mathrm{LSF}}
\newcommand{\mysub}[1]{( #1 )^\wedge}
\newcommand{\harm}{G}
\newcommand{\rmPV}{\PV}
\newcommand{\jac}{\Jac}
%\renewcommand{\sharp}{\#}
\newcommand{\myspace}{\mathsf}

\newcommand{\mes}{\delta_\Gamma}


\newcommand{\hs}{\raisebox{1pt}{$\scriptstyle \bigstar$}}

\newcommand{\scrb}{{\scriptscriptstyle \bullet}}

\newcommand{\rmins}[1]{ {\textrm{ #1 } } }

\newcommand{\et}{ \textrm{ et } }
\newcommand{\ou}{ \textrm{ ou } }
\newcommand{\avec}{ \textrm{ avec } }

%\newcommand{\og}{ \textrm{ og } }
%\newcommand{\eller}{ \textrm{  eller } }
\newcommand{\og}{ \ \land \  }
\newcommand{\eller}{ \ \lor \  }
\newcommand{\ikke}{ \ \urcorner }

\newcommand{\med}{\textrm{ med }}


\newcommand{\with}{\textrm{ with }}
\newcommand{\for}{\textrm{ for }}

\newcommand{\myand}{\textrm{ and }}
\newcommand{\myor}{\textrm{ or }}

\newcommand{\myif}{\textrm{ if }}
\newcommand{\myin}{\textrm{ in }}

\newcommand{\mywith}{\with}
\newcommand{\myfor}{\for}

\newcommand{\first}[1]{#1}


\newcommand{\infsup}[4]{\inf_{#1}\sup_{#2}\frac{#3}{#4}}
\newcommand{\cleq}{\preccurlyeq}
\newcommand{\cgeq}{\succcurlyeq}
\newcommand{\ceq}{\approx}

%\newcommand{\divides}{\ \backslash \! \backslash \ }
\newcommand{\divides}{\,  |  \, }

\newcommand{\whitney}{\mathfrak{W}}


%\newcommand{\orient}{\mathrm{o}}
\newcommand{\diff}{\mathrm{d}}
\newcommand{\rest}{\mathrm{r}}
\newcommand{\eval}{\mathrm{e}}
%\newcommand{\const}{\mathrm{c}}
\newcommand{\trans}{\mathrm{t}}
\newcommand{\curv}{\mathrm{c}}
\newcommand{\forget}{\mathrm{f}}

%\renewcommand{\orient}{\mathsf{o}}
\DeclareMathOperator{\orient}{\mathrm{o}}
\DeclareMathOperator{\eder}{\rmd}
\DeclareMathOperator{\pull}{\mathrm{pu}}
\DeclareMathOperator{\trace}{\mathrm{tr}}

\DeclareMathOperator{\clo}{\mathrm{clo}}
\DeclareMathOperator{\sta}{\mathrm{sta}}

\DeclareMathOperator{\ext}{\mathsf{ext}}

\renewcommand{\diff}{\mathsf{d}}
\renewcommand{\rest}{\mathsf{r}}
\renewcommand{\eval}{\mathsf{e}}
\renewcommand{\trans}{\mathsf{t}}
%\renewcommand{\const}{\mathsf{c}}
\renewcommand{\curv}{\mathsf{c}}
\renewcommand{\forget}{\mathsf{f}}

\DeclareMathOperator{\poincare}{\frP}
%\DeclareMathOperator{\poin}{\frp}
\DeclareMathOperator{\koszul}{\frK}


\DeclareMathOperator{\lie}{\calL}
\DeclareMathOperator{\jet}{jet}

%\renewcommand{\pull}{\mathrm{pb}}
%\renewcommand{\trace}{\mathrm{tr}}


%\newcommand{\transp}{\mathrm{T}}


\newcommand{\ula}[1]{\underleftarrow{#1}}


\newcommand{\deltat}{\delta'} %{\delta^\textsc{T}}
\newcommand{\deltal}{\tau}

\newcommand{\subcell}{\unlhd} % {\vartriangleleft}
\newcommand{\subcellstrict}{\lhd}
\newcommand{\subcelln}{\lhd^\cdot }
\newcommand{\subcellp}{\lhd_\cdot }
%\newcommand{\subcells}{\calS}
%\newcommand{\subcells}[1]{\overline{#1}}
\DeclareMathOperator{\subcells}{Cl}


\newcommand{\llp}{( \! (}
\newcommand{\rrp}{) \! )}
\newcommand{\lls}{[ \! [}
\newcommand{\rrs}{] \! ]}
\newcommand{\llc}{\{ \! \{}
\newcommand{\rrc}{\} \! \}}

\newcommand{\lsb}{ {\Big [} }
\newcommand{\rsb}{ {\Big ]} }

\newcommand{\ljump}{\lls}
\newcommand{\rjump}{\rrs}

\newcommand{\ctr}{ \, \mathsf{L} \, }

\newcommand{\Pe}{\textrm{P\'e}}

%\newcommand{\poly}{\bbP}
%\newcommand{\alter}{\bbA}
%\newcommand{\sym}{\bbS}
%\newcommand{\lin}{\bbM}

%\newcommand{\poly}{\calP}

\newcommand{\poly}{\rmP}
\newcommand{\alter}{\mathrm{Alt}}
\newcommand{\Sym}{\mathrm{Sym}}
\newcommand{\sym}{\mathrm{sym}}
\newcommand{\lin}{\mathrm{Lin}}
\newcommand{\Poly}{\calP}
\newcommand{\Pminus}{\Poly_{\!\!-}}

\newcommand{\Whitney}{\rmW}



\newcommand{\cochain}{\calC}
\newcommand{\bary}{\calB}
\newcommand{\refi}{\calR}


\newcommand{\cont}{\rmC}
\newcommand{\leb}{\rmL}
\newcommand{\sob}{\rmH}

\newcommand{\difo}{\Lambda}

\newcommand{\test}{\prime}

%-------------------math assembled------ 




\newcommand{\beq}{\begin{equation}}
\newcommand{\eeq}{\end{equation}}



\newcommand{\predicate}[2]{ \langle\! \langle\, #1  \,  :  \,  #2 \, \rangle \! \rangle} 

\newcommand{\setof}[2]{\{ #1 \ : \  #2 \}}
 
\newcommand{\twovec}[2]{
\begin{bmatrix}
#1\\
#2
\end{bmatrix}
}

\newcommand{\twomat}[4]{
\begin{bmatrix}
#1& #2 \\
#3& #4
\end{bmatrix}
}


\newcommand{\mapping}[4]{
\left\{
\begin{array}{rcl}
\displaystyle #1  &\to& #2\\
\displaystyle #3  &\mapsto & #4
\end{array} \right.
}

\newcommand{\threecases}[6]{
\left\{
\begin{array}{ll}
 #1 \ & \textrm{ if } #2 \\
 #3 \ & \textrm{ if } #4 \\
 #5 \ & \textrm{ if } #6 
\end{array}
\right.
}
\newcommand{\twocases}[4]{
\left\{
\begin{array}{ll}
 #1 \ & \textrm{ if } #2 \\
 #3 \ & \textrm{ if } #4
\end{array}
\right.
}
\newcommand{\totilfeller}[4]{
\left\{
\begin{array}{ll}
 #1 \ & \textrm{ hvis } #2 \\
 #3 \ & \textrm{ hvis } #4
\end{array}
\right.
}

\newcommand{\mixedleft}[4]{
\left\{
\begin{array}{l}
#1  \in #2 \\
#3  \in #4
\end{array}
\right.
}

\newcommand{\mixedright}[8]{
\
\left\{
\begin{array}{llll}
\forall #1  \in #2 \ & #3  & = #4\\
\forall #5  \in #6 \ & #7  & = #8
\end{array}
\right.
}










%--------------colors---------------------------

\definecolor{myyellow}{rgb}{1, 0.85, 0.15}
\definecolor{myorange}{rgb}{0.85, 0.5, 0.15}



\definecolor{mymagenta}{rgb}{0.85, 0.15, 1}
\definecolor{mycyan}{rgb}{0.15, 1, 0.85}


\definecolor{myred}{rgb}{0.7, 0.15, 0.15}
\definecolor{mypurple}{rgb}{0.425, 0.15, 0.425}
\definecolor{myblue}{rgb}{0.15, 0.15, 0.7}
\definecolor{mybluegreen}{rgb}{0.15, 0.425, 0.425}
\definecolor{mygreen}{rgb}{0.15, 0.7, 0.15}


\definecolor{mylightestbrown}{rgb}{0.825, 0.45, 0.225}
\definecolor{mybrown}{rgb}{0.55, 0.3, 0.15}
\definecolor{mydarkbrown}{rgb}{0.275, 0.15, 0.075}




\definecolor{mylightestgray}{rgb}{0.667,0.667,0.667}
\definecolor{mylightgray}{rgb}{0.5,0.5,0.5}
\definecolor{mygray}{rgb}{0.333, 0.333, 0.333}
\definecolor{mydarkgray}{rgb}{0.167,0.167,0.167}


%\FPeval{\colorvar}{round(1/7, 3)}

%\definecolor{myyellow}{rgb}{1, \real{1-\colorvar}, \real{\colorvar}}
%\definecolor{myorange}{rgb}{0.85, 0.5, 0.15}

%\definecolor{mypurple}{rgb}{0.425, 0.15, 0.425}

%\definecolor{myred}{rgb}{1-2*\colorvar, \colorvar, \colorvar}
%\definecolor{myblue}{rgb}{0.15, 0.15, 0.7}
%\definecolor{mygreen}{rgb}{0.15, 0.7, 0.15}

%\definecolor{mybrown}{rgb}{0.55, 0.3, 0.15}
%\definecolor{mydarkbrown}{rgb}{0.275, 0.15, 0.075}

%\definecolor{mylightgray}{rgb}{0.5,0.5,0.5}
%\definecolor{mygray}{rgb}{0.333, 0.333, 0.333}
%\definecolor{mydarkgray}{rgb}{0.167,0.167,0.167}






\newcommand{\red}[1]{{\textcolor{myred}{#1}}}
\newcommand{\green}[1]{{\textcolor{mygreen}{#1}}}
\newcommand{\blue}[1]{{\textcolor{myblue}{#1}}}

\newcommand{\gray}[1]{{\textcolor{mygray}{#1}}}
\newcommand{\brown}[1]{{\textcolor{mybrown}{#1}}}
\newcommand{\orange}[1]{{\textcolor{myorange}{#1}}}

\newcommand{\redbullet}{{\red{$\bullet$ }}}

\newcommand{\trm}[1]{
\textrm{#1}
}

\newcommand{\eqcom}[1]{
\quad \textrm{#1} }

\newcommand{\todo}[1]{
\red{\textbf{[#1]}}
}


\newcommand{\dikttittel}[1]{\begin{center}
\ttfamily \Large
#1
\end{center}
\vspace{1cm}
}


%------------------points------



\newcounter{npoint}[section]
\newcommand{\npoint}{
\paragraph{(\roman{npoint})}
\stepcounter{npoint}
}

\newcommand{\bpoint}{
\paragraph{$\bullet$ }
}

\newcommand{\ppoint}{
\paragraph{\pilcrow}$\!\!\!$
}

\newcommand{\case}[1]{
\bigskip
\noindent -- \emph{#1}
}

%\newcommand{\describe}[1]{{\bf (#1)}}

\newcommand{\myitem}[1]{\item {\bfseries \sffamily #1. }}
\newcommand{\myiteme}[1]{\item \emph{#1.} }

\newcommand{\remember}[1]{\blue{#1}}

\newcommand{\says}[1]{\vspace{1em}
{\sc #1} --
}

\newcommand{\mynewline}{\mbox{}\\}



%---------------header---------------
\newcommand{\myheading}[2]{
\begin{center}
{\Large \bf Notes for #1 -- #2\\
 Snorre Christiansen, \today\\
}
\rule{\textwidth}{1pt}\\
\vspace{1cm}
\end{center}
}



%-----------theorems--------------
\newtheorem{theorem}{Theorem}[section]
\newtheorem{lemma}[theorem]{Lemma}
\newtheorem{corollary}[theorem]{Corollary}
\newtheorem{proposition}[theorem]{Proposition}
\newtheorem{conjecture}[theorem]{Conjecture}

%\newtheorem{remark}{Remark}[section]

\theoremstyle{remark}
%\newtheorem{remark}{Remark}[section]
\newtheorem{example}{Example}[section]
\newtheorem{exercise}{Exercise}[section]
\newtheorem{problem}{Problem}[section]



\theoremstyle{definition}
\newtheorem{remark}{Remark}[section]
\newtheorem*{unremark}{Remark}
\newtheorem{definition}{Definition}[section]
\newtheorem{axiom}{Axiom}[section]


\newcounter{quest}[exercise]
\newcommand{\question}{
\medskip
\stepcounter{quest}
\noindent (\roman{quest}) 
}

%\newcommand{\hint}[1]{Hint in footnote\footnote{\emph{Hint:} #1}.}
\newcounter{subquest}[quest]
\newcommand{\subquestion}{
\smallskip
\stepcounter{subquest}
\noindent (\roman{quest}.\alph{subquest}) 
}


\newcommand{\solution}[1]{
%\begin{tcolorbox}[breakable, title= Suggested solution]
%#1
%\end{tcolorbox}
}

%\newcommand{\hint}[1]{
%  Hint in footnote\footnote{#1}.
%}


%--------------proofs-------------------------

%\newenvironment{proof}
%{\noindent{\itshape -- Proof: }}
%{\hfill $\square$}


\newcommand{\pushlength}{1cm}
\newcounter{pushlevel}[theorem]

\newcommand{\push}{
\stepcounter{pushlevel}
 \vspace{-1ex} \noindent \hspace{4.5ex} \begin{minipage}[t]{\textwidth*\real{0.99} - 4.5ex}
\mbox{}\hspace{-1ex}\rule[-1.5ex]{.1pt}{1.5ex}\rule{5.3ex}{.1pt}
\vspace{-.65ex}

}

\newcommand{\unpush}{
\end{minipage}
\noindent \mbox{}\hspace{4.2ex}\rule{.1pt}{1.5ex}\rule{5.3ex}{.1pt} \hspace{\stretch{1}}
\noindent
%\hspace{4ex} \rule{\textwidth - 4ex}{.1pt}
\addtocounter{pushlevel}{-1}

}


\newenvironment{proofof}[1]
{\noindent {\itshape -- Proof of #1:} }
{\hfill $\square$}


\newcounter{theopoint}[theorem]
%\newcounter{theopoint}[teorem]
\newcommand{\tpoint}{
\medskip
\stepcounter{theopoint}
\noindent (\roman{theopoint}) 
}




%-----------francais------------
\newtheorem{theoreme}{Th\'eor\`eme}[section]
%\newtheorem{proposition}[theoreme]{Proposition}
\newtheorem{lemme}[theoreme]{Lemme}
\newtheorem{corollaire}[theoreme]{Corollaire}








%---------------------quote--------------------------------



\newcommand{\myquote}[3]
{
\medskip
\begin{flushright}
\begin{minipage}{10.5cm} 
 #3
\end{minipage}\\
\bigskip
{\scshape #1} \  {\itshape #2}
\end{flushright}
\medskip
}

\newcommand{\tquote}[1]
{<<#1>>
}

\newenvironment{smallprint}
{
\bigskip
\begin{flushright}
\begin{minipage}{\textwidth*\real{0.833}} 
\footnotesize
}{
\normalsize
\end{minipage}
\end{flushright}
\bigskip
}





%------------------bibliography------
%\newcommand{\bibsep}[1]{\rule{5cm}{2pt}\\{\sffamily \bfseries #1}}

\newcommand{\bibref}[1]{[\ref{#1}]}
\newcommand{\biblabel}[1]{\label{#1}}
\newcommand{\bibauthor}[1]{{\scshape #1}: }
\newcommand{\bibtitle}[1]{{\itshape #1}; }
\newcommand{\bibspec}[1]{{\upshape #1}.}
\newcommand{\mybib}[3]{\bibauthor{#1}\bibtitle{#2}\bibspec{#3}}












%\input{../Definitions/definitionsSym}
%\input{../Definitions/definitionsEnv}

\numberwithin{equation}{section}

\title{Interpretation of a Discrete de Rham method % and\\ mixed Virtual Element Methods
  \\ as a Finite Element System}


\author{Snorre H. Christiansen\thanks{Department of Mathematics, University of Oslo, PO Box 1053 Blindern, NO 0316 Oslo, Norway. email: {\tt snorrec@math.uio.no.}}, Francesca Rapetti\thanks{Université Côte d'Azur, Laboratoire de Mathématiques J.A. Dieudonné, Parc Valrose, F-06108 NICE Cedex 02, France.}}

\date{}

\begin{document}

\maketitle

\begin{abstract}
  We show that the DDR method can be interpreted as defining a computable consistent discrete $\rmL^2$ product on a conforming FES defined by PDEs.  Without modifying the numerical method itself, this point of view provides an alternative approach to the analysis. The conformity and consistency properties we obtain are stronger than those previously shown, even in low dimensions. We can also recover some of the other results that have been proved about DDR, from those that have already been proved, in principle, in the general context of FES. We also bring VEM, the Virtual Element Method, into the discussion.
\end{abstract}

\bigskip

\noindent {\bfseries MSC:} 65N30, 58A12.

\bigskip

\noindent {\bfseries Keywords:} finite elements, differential forms.

\bigskip

\section*{Introduction}
We give a new presentation of the DDR method of \cite{DiPDroRap20, DiPDro23, BonDiPDroHu25}. We interpret it as a numerical method based on discrete spaces defined piecewise by solving certain new Partial Differential Equations (PDEs), on each cell of a mesh. The method is conforming in the sense of producing subcomplexes of natural complexes of Sobolev spaces. Even though the solutions to these new local PDEs are not computable, the numerical method itself is ''fully discrete'' because it allows the explicit computation of enough moments against polynomials, from the degrees of freedom. The method is then consistent in the sense of Strang's lemmas.  This type of consistency and conformity are stronger than those previously proved for DDR. The interpretation also embeds the method in the Finite Element System framework, and allows to leverage the results already obtained in that setting. 

\paragraph{Previous work on FES} The framework of Finite Element Systems (FES) was introduced in \cite{Chr08M3AS} and expanded upon in \cite[\S 5]{ChrMunOwr11} (this independent section is available on arXiv\footnote{\textsc{S. H. Christiansen,} \emph{Finite element systems of differential forms,} arXiv:1006.4779, 2015.}) and subsequent papers. It is an environment in which to define and analyze sequences of spaces of differential forms, that are finite elements in a broad sense. The underlying mesh can consist of cells (in particular polytopes) and the basis functions can be general (not necessarily based on polynomials). Thus one may call them polytopal discrete de Rham sequences, though this expression is given an a priori different meaning below. These spaces were conforming in the sense that the fields were $\rmL^2$-integrable with exterior derivative also $\rmL^2$-integrable, which enforces that traces, in the sence of pullbacks to faces, coincide (they are ''single-valued''). This encodes a type of partial continuity that is common for the mixed finite elements of electromagnetics and fluid flow \cite{BofBreFor13}, such as those of Raviart-Thomas-Nédélec. The latter were extended to differential forms in \cite{Hip99}. A synthesis with extensions to elasticity is provided in Finite Element Exterior Calculus (FEEC) \cite{ArnFalWin06}.

Since then the framework of FES has been generalized to accomodate higher regularity of the fields \cite{ChrHu18} (as in so-called ''Stokes complexes'' relevant to fluid flow) and other types of fields \cite{ChrHu23} (e.g. symmetric tensors and their associated curvature tensors, as in Riemannian geometry). The definition has a natural expression in the language of categories: a FES is a just a contravariant functor from the (poset) category determined by the mesh, to the category of vector spaces (interpreted as spaces of fields on cells). A short summary is provided in \cite{Chr22}. This definition covers all finite elements described according to Ciarlet's definition, including non-conforming ones (such as Morley and Crouzeix-Raviart elements).

This point of view connects finite elements with sheaf theory \cite{Ive86}. Indeed the combinatorial structure of the mesh can be seen as defining a ''site'' with a Grothendieck topology. A contravariant functor from such a site into an algebraic category is a presheaf, which for us represents function spaces than can be both localized (from global to local) and glued back (from local to global).

For the purposes of this paper, however, the definition given in \cite{Chr08M3AS} is sufficient. The goal back then was to define a discrete de Rham finite element sequence on polytopal meshes that mimicked the lowest order mixed finite elements that correspond to Whitney forms. Indeed these were known only for simplices (and products of simplices handled by tensor product constructions). For this presentation we will also refer to the more comprehensive presentation in \cite[\S 5]{ChrMunOwr11}. Short presentations can also be found in later publications, in particular \cite[\S 2]{ChrRap16}.

A high order version of the construction of \cite{Chr08M3AS} yielding minimal spaces (sometimes referred to as ''serendipity spaces'') was given in \cite{Chr10CRAS, ChrGil16}. The method of \cite{Chr08M3AS, ChrGil16} was conceptually based on harmonic extensions of differential forms, which is not in general computable. But with the help of, for instance, mesh refinements and known discrete spaces on the refinement, adequate discrete harmonic extensions were defined, yielding ''fully discrete'' numerical methods with algebraic and analytical properties identical with those of standard mixed finite elements.

By modifying the harmonic extension PDE that underlied the above constructions, one can obtain upwinded complexes, to treat convection dominated diffusion problems \cite{Chr13FoCM,ChrHalSor14}. The framework has also been used to give compact descriptions of several known finite element methods \cite{ChrGil16, ChrRap16}. The generalized Whitney forms introduced in \cite{Chr08M3AS} have also been used to analyze Discrete Exterior Calculus (DEC) in \cite{GuzPot25}.

\paragraph{Other polytopal methods} The Virtual Element Method (VEM), reviewed in \cite{BeiBreMarRus23}, has been developed with similar goals as FES and some similar basic principles, such as definition by recursive extension of the fields, from low dimensional cells. But, crucially, no submesh is used in VEM to define the discretization spaces or the numerical method. The challenge is to reconstruct enough information on the fields from the degrees of freedom, to obtain stable and consistent approximations of the bilinear forms defining the partial differential equation to solve. A connection between the mixed VEM (defined in dimension 3) of \cite{BeiBreMarRus16} and FES (in arbitrary dimension) was made in \cite[\S 2.1]{ChrGil16} and will be used later.  

The Discrete de Rham (DDR) method was introduced in dimension 2 and 3 in \cite{DiPDroRap20} and has since been further developed in several papers. We refer to \cite{BonDiPDroHu25} for an overview and a description of the method valid for differential forms of arbitrary degree in arbitrary dimension. 

The DDR construction was inspired at least in part by Hybrid High Order (HHO) methods. Accordingly, it has been presented as a non-conforming method involving spaces of non-matching fields on cells of all dimension. Non-conformity of the spaces poses major challenges for the analysis of the method, many of which have been addressed in the literature \cite{DiPDro23,DiPDroPit23}. In particular an original notion of ''adjoint consistency'' has been developed.

The close relationship between DDR and VEM has already been pointed out in \cite{BeiDasDiPDro22}. In some versions these methods can be seen to have the same degrees of freedom, and they coincide with those of FES provided in \cite[\S 2.1]{ChrGil16}.

Lowest order numerical methods have been designed on various polytopes for a long time and can be traced back to work of Yee and Arakawa in the 1960ies and 1970ies. Many seem to share a common algebraic approach to the differential operators whereas they have different approaches to the construction of the discrete $\rmL^2$ inner product.

\paragraph{Contribution} In this paper, we reinterpret DDR as a conforming FES equip\-ped with computable consistent approximations of the $\rmL^2$ inner products.  The degrees of freedom of \cite[\S 2.1]{ChrGil16} are used.  This shows that these examples of DDR, VEM and FES are very similar, if not identical, in practice. They seem to differ mainly in terms of interpretation and analysis. The main novelty here is just that a new harmonic extension PDE, modifying the one used in \cite[\S 2.1]{ChrGil16}, is introduced to define the discrete spaces.

This point of view allows us to recover some of the results that have been proved for VEM or DDR %\cite{DiPDro23,DiPDroPit23},
as special cases of results already proved, in their principle, in the general context of FES.  We summarize:
\begin{itemize}
\item Consistency of the discrete $\rmL^2$ scalar product now follows from the fact that the new local PDE allows the explicit recursive computation of $\rmL^2$ products with enough polynomials, from the degrees of freedom.
\item Cohomological properties developed for FES, especially \cite[Proposition 5.15]{ChrMunOwr11}, now carry over to DDR. This passage from local exactness (on each cell of the mesh), to globally correct cohomology (on the computational domain) illustrates that the existence of unisolvent degrees of freedom on a FES can be interpreted as a softness condition on sheaves (see \cite[\S III.2]{Ive86} for definitions).
\item Uniformly bounded commuting projections, building on \cite{Sch08, Chr07NM}, have been described for FEEC \cite{ArnFalWin06,ChrWin08} and FES \cite[\S 5.3]{ChrMunOwr11}. They enable one to recover Poincaré-Friedrich estimates, which are basic to the stability analysis of mixed methods.
\end{itemize}
In the setting of FES one furthermore obtains refinements in the analysis:
\begin{itemize}
\item Translation estimates are given in \cite[Proposition 5.71]{ChrMunOwr11}, and have applications to the equations of fluid flow.
\item Discrete Rellich compactness, as introduced by Kikuchi for Maxwell \cite[\S IV.19]{Bof10}, holds generally, as proved first in \cite[Corollary 5]{Chr07NM} with applications to the Hodge-Laplace eigenvalue problem on manifolds. See also \cite{ChrWin13IMA} on the relationship of bounded commuting projections to such properties, at an abstract level.
\item Discrete Sobolev injections hold \cite[Propositions 5.69, 5.70]{ChrMunOwr11} and have applications to, for instance, non-linear Maxwell wave equations as in \cite{ChrSch11}. See also \cite{HeHuXu19} on this topic.
\item Fractional order Sobolev estimates are provided in \cite{Chr18}, with applications to the Dirac equation.
\end{itemize}

We insist that in our point of view the method is conforming in the sense that we obtain subspaces of standard Sobolev spaces and that the discrete differential operators (instances of the exterior derivative) are simply the restriction of the ordinary ones. These operators are computable in terms of the chosen degrees of freedom. Taking into account the variational crime arizing from the approximation the $\rmL^2$ inner product is viewed as a separate ingredient in the analysis, addressed with the help of Strang's first lemma.

An example of a FES determined by solving local PDEs and equipped with a computable and consistent bilinear form, has been analyzed in \cite{ChrHal15}. It concerned the Schrödinger equation with a magnetic field (representing an $\bbU(1)$-connection). The local PDEs used on the cells are often related, but not identical, to the global PDE one wants to solve. This aspect can be compared with Trefftz methods which often glue exact local solutions in non-conforming ways \cite{HipMoiPer16}.

Notice that many arguments involve a straightforward scaling combined with an appeal to compactness inspired by \cite[Remark p. 64 and proof p. 69]{ArnFalWin06}. An optimal formulation of such compactness arguments remains to be delineated in the context of FES.

\begin{unremark}
More recent constructions of \emph{local} bounded commuting projections \cite{FalWin14, ArnGuz21} carry over to FES in a very natural way. The approach of \cite{FalWin14} extends to FES of differential forms and the approach of \cite{ArnGuz21} extends to general FES as in \cite{ChrHu23} (the latter approach does not appeal to a wedge product). In fact the FES framework seems to identify exactly the algebraic properties of discretization spaces needed for these constructions. A general analysis of stability and compactness properties of these methods remains, in the context of FES, to be fleshed out. See \cite{ErnGuzPotVoh25} for the state of the art.
\end{unremark}

In a sense we have little to add to the design of the DDR and VEM methods themselves. Many of the defining integrations (by parts or not) against cleverly chosen polynomial forms have of course already been exemplified in the literature, though perhaps not in this order and generality. We hope the FES perspective can contribute usefully to the analysis, leading to various convergence proofs. Besides, by providing a clarified rationale for neglecting certain terms and identifying key underlying algebraic facts (such as short exact sequences of complexes and their long exact sequences of cohomology groups), this presentation can ease the extension of the method in various directions. This could be for instance the above mentioned upwinding problem, or the existing variants of DDR developed for other differential operators \cite{DiPHan24} and regularities \cite{Han23}.

The paper is organized as follows. In section \ref{sec:not} we review notations and provide a reading guide on FES. Section \ref{sec:low} contains the construction in the lowest order case (sequences of spaces containing constants). Section \ref{sec:high} contains the construction in the high order case (sequences of spaces containing polynomials of degree $r$, for any choice of $r \geq 1$). Section \ref{sec:algo} shows how the preceding results can be used to set up a numerical method for a prototype equation and start to analyse it.

\section{Notations\label{sec:not}}
 Let $\calT$ be a polytopal mesh of some Euclidean polyhedral domain. When considering a polytope $T$ with a face $T'$, if $u$ is a differential form on $T$, by its trace on $T'$ we always mean the trace in the sense of pullback of $u$ to $T'$ by the inclusion of $T'$ in $T$.

For each $k \in \bbN$ and $T\in \calT$ we let $X^k(T)$ be the space of differential $k$-forms on $T$ which are in $\rmL^2(T)$ with exterior derivative in $\rmL^2(T)$. The $\rmL^2$ norm is denoted $|\cdot|$ and the natural graph norm on $X^k(T)$ is denoted $\| \cdot \|$.


The subspace of $X^k(T)$ consisting of forms with $0$ traces on $\partial T$ is denoted $X^k_0(T)$:
\begin{equation}
  X^k_0(T) = \{ u \in X^k(T) \ : \ u|_{\partial T} = 0\}.
\end{equation}
The subscript $0$ will have the same meaning for other spaces as well. Elements on $X^k(T)$ have well defined traces, in negative index Sobolev spaces, only on codimension one faces $T' \subseteq \partial T$. However, if these traces in turn are imposed to be in $X^k(T')$, one can repeat the procedure and follow traces all the way down. One gets a space of forms with improved boundary regularity and some nice properties, such as \cite[Proposition 2.4]{ChrRap16}. All our discrete spaces will have this type of extra boundary regularity.

 % Furthermore, the pullback to any face $T'$ is also in $\rmL^2(T')$ with exterior defivative in $\rmL^2(T')$. Notice that this makes sense because such traces are well-defined a priori on codimension one faces, in Sobolev spaces of negative index.

The space of Whitney $k$-forms (which could also be attributed to Weil and de Rham) was originally presented in terms of partitions of unity, such as barycentric coordinates, but it can also be presented in terms of the Koszul operator \cite[\S 3.2]{ArnFalWin06}. For each cell $T \in \calT$, we let $b_T$ be its barycentre (we suppose it is ''computable''). On $T$, the Koszul operator $\kappa_T$ is the contraction by the vector field on $T$ given as $x \mapsto x- b_T$ (for $x \in T$). We adopt notations from \cite{ArnFalWin06}. The original lowest order Whitney forms are then $\Pminus^1\Lambda^k(T)$. The high order version are denoted $\Pminus^r\Lambda^k(T)$:
\begin{equation}
\Pminus^r\Lambda^k(T) = \{ u \in \Poly^r \Lambda^k(T) \ : \ \kappa_T u \in \Poly^r \Lambda^{k-1} (T)\}.
\end{equation}
These spaces were defined in \cite{Hip99} in terms of the standard Poincaré null-homotopy operators rather than the Koszul operators and were described in FES terms in \cite{ChrRap16} under the name ''trimmed'' polynomial differential forms. 


The Hodge star operator is denoted $\hs$. It maps from $k$-forms on $T$ to $(d-k)$-forms on $T$, when $d = \dim T$. It is characterized by the property that, when $u$ and $v$ are both $k$-forms:
\begin{equation}
  u \wedge \hs v = (u \cdot v)\, \omega,
\end{equation}
where the pointwise scalar product is denoted $(\cdot)$ and the volume form on $T$ associated with the scalar product and the orientation is denoted $\omega$. For instance $\hs 1 = \omega$.
The formal adjoint of $\rmd$ with respect to the $\rmL^2$ scalar product is denoted $\rmd^\star$. We have $\rmd^\star = \pm \hs \rmd \hs$.

A third point of view on Whitney forms, developed in \cite{Chr08M3AS}, is that they are the solutions to the PDE system:
\begin{equation}
  \rmd^\star \rmd u = 0, \quad \rmd^\star u = 0,
\end{equation}
supposed to hold on simplices of all dimensions (equipped with translation invariant metrics). For $0$-forms this reduces to $\Delta u = 0$. For $k$-forms on simplices (or flat cells for that matter) of dimension $k$, this reduces to $\rmd^\star u= 0$, which is a characterization of the constant $k$-forms.

Even though we adopt the FES point of view, we first present our results without relying on those definitions. On the other hand \cite[\S 2]{ChrRap16} contains a glossary and a short presentation that is more than sufficient for this article.

\paragraph{FES: an overview}  Since the FES framework has not been reviewed recently we now offer a guide to the relevant literature, for the convenience it might provide the reader.
\begin{itemize}
\item In \cite{Chr08M3AS} the FES framework was introduced as auxiliary discrete spaces inside which analogues of (lowest order) Whitney forms could be defined as locally harmonic forms.
\item Reference \cite{Chr09AWM} is mostly expository.  In \S 3.1 some background in algebra is provided (the five lemma, the snake lemma and the long exact cohomology sequence, the Künneth theorem on tensor products). In \S 3.2 basic concepts in differential geometry are reviewed. This algebraic and differential geometric background is then applied to study FES. The motivation was the construction of Galerkin methods for Maxwell's equations, complementing \cite{Hip02,Jol03,Mon03} on some points, particularly in regards to general meshes and basis functions.

  A ''geometric decomposition'' in terms of ''bubble spaces'', in the sense of \cite{ArnFalWin09}, is given in \S 3.3.2 (Proposition 3.14) as a byproduct of discretely harmonic forms.

  A first study of tensor products of FES is provided in \S 3.3.3 and further developed in later publications. This theory covers standard elements as in \cite[\S 3]{ArnBofBon15} and splines as in \cite{BufRivSanVaz11}. In fact, one can start with any space of functions on an interval that contains a non-trivial constant function and one with different values at the endpoints. The latter function can be for instance an exponential. 

\item In \cite[\S 5]{ChrMunOwr11} the framework was further developed in its own right, to account for the mixed finite element spaces synthesized in FEEC. In \S 5.3 commuting quasi-interpolators deduced from smoothings were studied. In \S 5.4 stability properties were deduced.

\item In \cite{Chr13FoCM} it is illustrated that by modifying the definition of harmonicity one can get upwinded complexes of differential forms, useful for convection diffusion problems. The above mentioned exponential function allows one to recover some tensor product ``exponentially fitted'' methods. The scalar convection diffusion problem is analyzed in \cite{ChrHalSor14} in fractional Sobolev spaces (of the type $\rmH^{1/2}$).
  
\item In \cite{ChrRap16} FES was used to describe the spaces $\Pminus^r\Lambda^k(T)$ of trimmed polynomial differential forms on simplices. In \S 2  a glossary and a seven page summary of FES is provided. Alternative degrees of freedom were proposed. Resolutions were introduced to describe relations in the natural spanning families constructed with barycentric coordinates.

  It is believed that working with canonical generators and relations (instead of selecting a specific basis) and systematically treating restrictions to subcells, could be valuable for implementations, not just the theoretical study. 

\item In \cite{ChrGil16} it was shown that the Serendipity element \cite{ArnAwa14}, the TNT element \cite{CocQiu14} and the trimmed polynomial differential forms are all examples of FES satisfying a minimality condition. A connection with VEM was made explicit in \S 2.1. A general recipy for constructing minimal or ''serendipity'' spaces was given in \S 3.

\item In \cite{ChrHu18} and \cite{ChrHu23} the framework of FES is generalized to account for more types of finite element spaces. In \cite{ChrHu23,Chr22} the connection with category theory and sheaves is described.
\end{itemize}


%%%%%%%%%%%%%%%%%%%%%%%%%%%%%%%%%%%%%%%%%%%%%%%%%%%%%%%%%%%%%%%%%%%%%%%%%%%%%%
\section{Lowest order case\label{sec:low}}
Let $T$ be a cell. The case $\dim T =0$ need not be detailed, so we suppose that $\dim T = d > 0$.

In this section we define $W^k(T)$ as a slight modification of $\Pminus^1 \Lambda^k(T)$, namely, that for, $k= 0$ we impose that the elements of $W^0(T)$ should have integral $0$.

Then we define:
\begin{align}
  Z^k_0(T) & = \{u \in X^k_0(T) \ : \ \forall w \in W^{d-k}(T) \quad \int_T u \wedge w = 0\}. % ,\\
%  Z^k_0(T) & = Z^k(T) \cap X^k_0(T).
\end{align}

\begin{proposition}\label{prop:zexact}
  The spaces $Z^\bs_0(T)$ form a sequence which is exact except at index $k=\dim T$, where the integral determines an isomorphism from the cohomology group to $\bbR$.
\end{proposition}
\begin{proof}
  We have a short exact sequence of complexes:
\begin{equation}
  0 \to Z^\bs_0(T) \to X^\bs_0(T) \to W^{d-\bs}(T)^\star \to 0.
\end{equation}
The sequence $W^\bs(T)$ is exact, including at indices $0$ and $d$. Hence $Z^\bs_0(T)$ inherits the cohomology of $X^\bs_0(T)$, by the snake lemma, in its full version. See \cite[Theorem 3.5, p. 600]{Alu09} or \cite[Theorem 3.1, p. 356]{Chr09AWM}.
\end{proof}


We define:
\begin{equation}
  Y^k(T) = \{ u \in Z^k_0(T) \ : \ u \perp \rmd Z^{k-1}_0(T) \}.
\end{equation}
Orthogonality is taken with respect to the $\rmL^2$ scalar product (also elsewhere in this paper).
We remark that Proposition \ref{prop:zexact} entails that the following is an isomorphism:
\begin{equation}
\rmd: Y^{k-1}(T) \to \ker \rmd |_{Z^{k}_0(T)},
\end{equation}
for $k < d$, and that, for the case $k =d$, we also have the isomorphism:
\begin{equation}
\rmd: Y^{d-1}(T) \to \ker \ts \int_T(\cdot) |_{Z^{d}(T)}.
\end{equation}


We define $B^k(T)$ to be the subspace of $X^k(T)$ consisting of fields $u \in X^k(T)$ such that the following system of PDEs holds:
\begin{align}
  \int_T \rmd u \cdot \rmd u' + \int_T \rmd u' \wedge \kappa_T a + \int_T u' \wedge \kappa_T b + \int_T u' \cdot \rmd v & = 0,\\
  \int_T \rmd u \wedge \kappa_T a' =0 ,\ \int_T u \wedge \kappa_T b' & = 0 ,\\
  \int_T u \cdot \rmd v' & = 0,
\end{align}
for all $u' \in X^k_0(T)$, for some $a \in \rmP^0\Lambda^{d-k}(T)$, some $b \in \rmP^0\Lambda^{d-k +1}(T)$, and some $v \in Y^{k-1}(T)$,  for all $v' \in Y^{k-1}(T)$, all $a' \in \rmP^0\Lambda^{d-k}(T)$ and all $b' \in \rmP^0\Lambda^{d-k +1}(T)$. 

Notice that $a', b'$ and $v'$ enforce constraints on $u$, and that $a, b, v$ can be interpreted as the corresponding Lagrange multipliers. We refer to $a'$ and $b'$ as the finite dimensional constraints.

\begin{remark}
Removing the finite dimensional space of constraints represented by $a,a'$ and $b,b'$, and allowing general $v,v'$, we obtain the harmonic extension used in \cite{Chr08M3AS}. Indeed the equations considered there were:
\begin{align}
  \rmd^\star \rmd u + \rmd v & = 0,\\
  \rmd^\star u & = 0.
\end{align}
The Lagrange multiplier $v$ was always $0$. This can be interpreted as solving a Hodge-Laplace problem. Indeed, with $-\Delta = \rmd^\star \rmd + \rmd \rmd^\star$, the above system is equivalent to:
\begin{equation}
  \Delta u = 0,\ \textrm{and }\ (\rmd^\star u) |_{\partial T} = 0.
\end{equation}
In both cases the boundary condition is that $u|_{\partial T}$ is given (in the sense of pullbacks).

In \cite{Chr08M3AS}, discrete variants obtained from weak formulations on suitable finite dimensional spaces, were also used to make the method ''fully discrete''.

To sum up the definition of $B^k (T)$ involves a variant of the Hodge-Laplace equation deduced for certain boundary conditions, modified by some polynomial constraints.
\end{remark}

We now detail the special cases $k=0$ and $k=d$ of the definition of $B^k(T)$.

\noindent For $k= 0$ we are left with the following system:
 \begin{align}
  \int_T \rmd u \cdot \rmd u' + \int_T \rmd u' \wedge \kappa_T a = 0,\\
  \int_T \rmd u \wedge \kappa_T a' = 0.
\end{align} 

\noindent For $k=d$, we are left with:
\begin{align}
  \int_T u' \wedge \kappa_T b + \int_T u' \cdot \rmd v = 0,\\
  \int_T u \wedge \kappa_T b' = 0 ,\\
  \int_T u \cdot \rmd v' = 0,
\end{align}
Notice that $\kappa_T b$ is a general element of $W^0(T)$. It has integral $0$ because of the choice of origin in $\kappa_T$. Let $\overline u$ be the average of $u$. Then $u-\overline u$ is in $Y^d(T)$ and has integral $0$, so it is $0$. That is $u = \overline u$. In other words $B^d(T)$ consists of the constant $d$-forms.

Then we define $A^k(T)$ to be the subspace of $X^k(T)$ consisting of fields $u$ whose trace on any face $T'$ of $T$ (including $T' = T$) is in $B^k(T')$.

We make the following claims, which we will justify subsequently:
\begin{itemize}
\item A field in $B^k(T)$, for $k < \dim T$ is uniquely determined by its boundary traces. Moreover the extension problem defined by the above PDE is wellposed.
\item The family of spaces $A^k(T)$ is closed under traces and exterior derivative, making them a FES.
\item The spaces $A^k(T)$ are finite dimensional and have one degree of freedom per $k$-face of $T$. This ensures \emph{softness} of the FES.
\item For each $T \in \calT$, the spaces $A^\bs(T)$ form an exact sequence resolving $\bbR$. Together with softness this ensures \emph{compatibility} of the FES.
\item The space $A^k(T)$ contains the Whitney forms $\Pminus^1 \Lambda^k(T)$. This ensures approximation properties.
\item For any $u \in A^k(T)$ and any $a \in \rmP^0\Lambda^{d-k}(T)$, the integral $\int_T u \wedge a$ can be computed exactly from the degrees of freedom, in a recursive way.
\end{itemize}


According to plan, we first comment on the well-posedness of the PDE. We will use the following.

\begin{proposition} On a cell $T$ of dimension $d$ and for $k < d$, the above system of PDEs for $B^k(T)$ is wellposed, for given boundary data.
\end{proposition}

\begin{proof}
Extend the boundary data to an arbitrary field $\tilde u \in X^k(T)$ (for instance by ordinary harmonic extension). Then write the unknown $u \in X^k(T)$ as $u = \tilde u + \hat u$ with $\hat u \in X^k_0(T)$. We now consider the system:
\begin{align}
  \int_T \rmd \hat u \cdot \rmd u' + \int_T \rmd u' \wedge \kappa_T a + \int_T u' \wedge \kappa_T b + \int_T u' \cdot \rmd v & = - \int_T \rmd \tilde u \cdot \rmd u',\\
  \int_T \rmd \hat u \wedge \kappa_T a' = - \int_T \rmd \tilde u \wedge \kappa_T a'  ,\ \int_T \hat u \wedge \kappa_T b' & = - \int_T \tilde u \wedge \kappa_T b' ,\\
  \int_T \hat u \cdot \rmd v' & =  - \int_T \tilde u \cdot \rmd v' ,
\end{align}


Since $k < d$ we can now integrate by parts the constraints, as in:
\begin{align}
\int_T \rmd \hat u \wedge \kappa_T a' + \int_T \hat u \wedge \kappa_T b' & = \int_T \hat u \wedge ( (-1)^{k+1} \rmd \kappa_T a' + \kappa_T b'),\\
&=  \int_T \hat u \wedge (\frac{(-1)^{k+1}}{d-k} a' + \kappa_T b') ,
\end{align}
For $k = 0$ there is no $b'$ term. %For $k = d$ there is no $a'$ term.
We recognize the general expression of a $w' \in \Whitney^{d-k}(T)$:
\begin{equation}\label{eq:wprime}
  w' = \frac{(-1)^{k+1}}{d-k} a' + \kappa_T b',
\end{equation}
uniquely corresponding to $a'$ and $b'$.

We can rewrite the system for  $u \in X^k_0(T)$ as, for all $u' \in X^k_0(T)$:
\begin{align}
  \int_T \rmd \hat u \cdot \rmd u' + \int_T u' \wedge w + \int_T u' \cdot \rmd v & = - \int_T \rmd \tilde u \cdot \rmd u',\\
  \int_T  \hat u \wedge w' & = - \int_T \rmd \tilde u \wedge \kappa_T a' - \int_T \tilde u \wedge \kappa_T b', \\
  \int_T \hat u \cdot \rmd v' & =  - \int_T \tilde u \cdot \rmd v' ,
\end{align}
for some $w \in W^{d-k}(T)$ and $v \in Y^{k-1}(T)$, and for all $w' \in W^{d-k}(T)$ (written as (\ref{eq:wprime}) ) and $v' \in Y^{k-1}(T)$.


We look first at the homogeneous problem (zero right sides). The second equation says that $\hat u \in Z^k_0(T)$ and the third one then says that $\hat u\in Y^k(T)$. The first equation with $u' = \hat u$ gives $\rmd \hat u =0$. This shows that $\hat u = 0$. The homogeneous problem thus has only the trivial solution.

The proof of the well-posedness in the PDE sense, is contained in the Lemmas \ref{lem:infsup} and \ref{lem:coerker}, below.
\end{proof}

 We begin with the following remark.
\begin{lemma} We have: $(\hs W^{d-k}(T)) \cap \rmd Y^{k-1}(T) = 0$.
\end{lemma}
\begin{proof}
 % Indeed we have $\hs W^{d-k}(T) \subseteq \poly^1 \Lambda^k(T)$.

 %  For $k< d$, such forms are determined by their boundary traces, which are zero for $\rmd Y^{k-1}(T)$.

It holds, more generally, that for $u \in Z^{k}_0(T)$ we have $u \perp \hs W^{d-k}(T)$ by definition.
\end{proof}


\begin{lemma}[Inf-Sup condition]\label{lem:infsup} Fix $k$ and $T$. Suppose that $E$ is a finite dimensional subspace of $X^k(T)$  and that $E \cap \rmd Y^{k-1}(T) = 0$. Then:
\begin{equation}
  \inf_{v \in \rmd Y^{k-1}(T), w \in E} \sup_{u\in X^k_0(T)} \frac{|\langle u, v \rangle + \langle u, w \rangle|}{\| u\| (\| v \| + \| w\|)} >0.
\end{equation}
\end{lemma}
\begin{proof}
  If not, construct sequences $v_n \in \rmd Y^{k-1}_0(T)$ and $w_n \in E$ such that $\|v_n\| +\|w_n\| = 1$ and:
  \begin{equation}
    \sup_{u\in X^k_0(T)} \frac{|\langle u, v_n \rangle + \langle u, w_n \rangle|}{\| u\|} \to 0.
  \end{equation}
  Extracting subsequences we suppose futhermore that $v_n \to v$ weakly in $X^k_0(T)$ and that $w_n \to w$ strongly in $E$ (by finite dimensionality of $E$).

  From the above we get that $v + w = 0$. Also $\rmd Y^{k-1}(T)$ is closed in $X^k_0(T)$, so $v\in \rmd Y^{k-1}(T)$. Therefore $v= 0$ and $w= 0$. Thus we conclude that $w_n \to 0$ strongly. From that we get:
  \begin{equation}
    \sup_{u\in X^k_0(T)} \frac{|\langle u, v_n \rangle|}{\| u\|} \to 0.
  \end{equation}
  But this is false, since we can use $u = v_n$ in this $\sup$.
  \end{proof}

\begin{lemma}[Coercivity on the kernel]\label{lem:coerker}
  There is $C>0$ such that for all $u \in Z^k_0(T)$, such that $u \perp \rmd Y^{k-1}(T)$ (w.r.t. $\rmL^2$ scalar product):
\begin{equation}
   |u| \leq C | \rmd u|.
\end{equation}
\end{lemma}
\begin{proof}
  If not choose a sequence $u_n \in Z^k_0(T)$, such that $u_n \perp \rmd Y^{k-1}(T)$, $|u_n| = 1$ and $|\rmd u_n| \to 0$. Extracting subsequences we suppose that $(u_n)$ converges weakly in $\rmL^2$ to some $u$. Then $u \in Z^k_0(T)$ and $\rmd u =0$ and $u \perp \rmd Y^{k-1}(T)$. Hence $u=0$.

  We now show that the convergence $u_n \to 0$ is actually strong in $\rmL^2$. For that purpose we control $\rmd^\star u_n$ in $\rmL^2$. We decompose, with respect to the $\rmL^2$ product:
  \begin{equation}
    X_0^{k-1}(T) = Z_0^{k-1}(T) \oplus F.
  \end{equation}
  Any $v\in X_0^{k-1}(T)$ such that $v \perp \rmd X^{k-2}_0(T)$, can now accordingly be written $v = v' + w$ with $v' \in Z^{k-1}_0(T)$ and $w \in F$. We have:
  \begin{equation}
   \sup_{v \in X_0^{k-1}(T)} \frac {|\int_T u_n \cdot \rmd v|}{| v|} \leq \sup_{w \in F} \frac{|\int_T u_n \cdot \rmd w|}{|w|}.
  \end{equation}
  Since $F$ is finite dimensional and $u_n\to 0$ weakly in $\rmL^2$, we get that $\rmd^\star u_n \to 0$ strongly in $\rmL^2$.

  Since the subspace of $X^k_0(T)$ of consisting of fields $u$ with $\rmd^\star u \in \rmL^2(T)$, with its natural topology, is compactly injected in $\rmL^2(T)$, $(u_n)$ converges strongly in $\rmL^2$ norm to $0$. But this contradicts $|u_n| = 1$.
  \end{proof}

\begin{proposition}[FES] The family of spaces $A^k(T)$ is closed under traces and the exterior derivative. 
\end{proposition}
\begin{proof}
  Closedness under traces follows from the definition of $A^k(T)$, that already took into account all the traces.

  Closedness under the exterior derivative will follow from the fact that if $u \in B^k(T)$ then $\rmd u \in B^{k+1}(T)$. Suppose then that $u \in B^k(T)$. We have for all $u' \in X^k_0(T)$:
  \begin{equation}
\int_T \rmd u \cdot \rmd u' + \int_T u' \wedge w + \int_T u' \cdot \rmd v = 0.
  \end{equation}
  If now $u'\in Y^k(T)$, the second and third terms are 0. This gives the last constraint on $\rmd u$.

The two finite dimensional constraints are trivially satisfied.
The top equation is also trivially satisfied, with zero Lagrange multipliers.
\end{proof}



\begin{proposition}[Degrees of Freedom]
On $A^k(T)$ the integrals on $k$-faces of $T$ constitute unisolvent degrees of freedom.
\end{proposition}
\begin{proof}
Any given $u \in A^k(T)$ is determined by its codimension one traces, that can be followed all the way down to faces of dimension $k$, where $u$ must be constant, hence determined by its integral. Conversely from any values of these integrals, a $u \in A^k(T)$ can be reconstructed by recursive extension of the corresponding constant forms on $k$-cells.
\end{proof}

\begin{proposition}[Local sequence exactness]
The sequence $A^\bs(T)$ is exact (more precisely it resolves $\bbR$).
  \end{proposition}
\begin{proof}
  The degrees of freedom provide a cochain isomorphism with the cellular cochain complex, which is known to be exact on any given cell.
\end{proof}

\begin{proposition}[Approximation]
  The spaces $A^k(T)$ contain $\Pminus^1 \Lambda^k(T)$.
\end{proposition}

\begin{proof}
Whitney forms satisfy $\rmd^\star u = 0$ and $\rmd^\star \rmd u = 0$, so the crucial points are the finite dimensional constraints.

  It was proved in \cite[Preprint, Proposition 3.9]{Chr07NM} that $\wedge : \Pminus^r \Lambda^k(T) \times \Pminus^s \Lambda^l(T) \to \Pminus^{r+s}\Lambda^{k+l}(T)$. A simpler proof based on the Koszul operator, rather than barycentric coordinates, was then provided in \cite[Equation (3.16)]{ArnFalWin06}. Based on this result, if $u\in \Pminus^1 \Lambda ^k(T)$ and $b' \in \rmP^0 \Lambda^{d-k+1}(T)$ we have that $u \wedge \kappa_T b$ is a linear $d$-form. Its integral can therefore be obtained from evaluation at the barycentre of $T$. Since furthermore $\kappa_T$ was constructed with respect to this point, we get that $\int_T u \wedge \kappa_T b' = 0$. Similarly $\int_T \rmd u \wedge \kappa_T a' = 0$ for $a' \in \rmP^0 \Lambda^{d-k}(T)$. 
\end{proof}

\begin{proposition}[Computation of moments]\label{prop:momconst}
For $u \in A^k(T)$, the integrals $\int_T u \wedge a'$ against constant $a' \in \rmP^0 \Lambda^{d-k} (T)$ can be computed recursively from the degrees of freedom.
\end{proposition}

\begin{proof}
  Let $T$ be a cell of dimension $d$. The case $k = d$ is immediate, so we suppose $k < d$.

  We integrate by parts the first constraint (as before, but now with a boundary term):
  \begin{equation}
    \int_T \rmd u \wedge \kappa_T a'  = \int_{\partial T} u \wedge \kappa_T a' + \int_T u \wedge \frac{(-1)^{k+1}}{d-k} a'. 
  \end{equation}
  The left side is zero by definition. On the right side, the second term is of the form that we would like to compute.

  We are left with the problem of evaluating the first term on the right side. It will be decomposed into integrals on the $(d-1)$ dimensional faces $T'$ of $T$. That is, we require terms of the form:
  \begin{equation}
    \int_{T'} u \wedge \kappa_T a'.
  \end{equation}
%In the case where $k = d-1$,  $\kappa_T a'$ has constant pullback to $T'$ and the above integrals are given directly by the degrees of freedom.
We use that $(\kappa_T -\kappa_{T'})a'$ is constant on $T'$ and that by definition:
  \begin{equation}
    \int_{T'} u \wedge \kappa_{T'} a' = 0.
  \end{equation}
The boundary term has been reduced to integrals of $u$ wedged against constants.

This method can be pursued, decreasing the dimension by one at each step, all the way down to the case where the integrals against constants are actually degrees of freedom. 
  \end{proof}

For a given $u \in A^k(T)$ one can then compute the $\rmL^2(T)$ projection of $u$ onto constants $\Poly^0 \Lambda^k(T)$, from its degrees of freedom. We denote it as $\gamma_T u$.  In practice one computes $\gamma_{T'} u$ starting with $\dim T' = k$ and increasing dimension of $T'$ by one at each step.

In \S \ref{sec:algo} we sketch how $\gamma_T$ is used as a crucial ingredient in the design of some numerical methods, that are amenable to an analysis based on Strang's first lemma \cite[Theorem 26.1 p. 192]{Cia91}.
  

%%%%%%%%%%%%%%%%%%%%%%%%%%%%%%%%%%%%%%%%%%%%%%%%%%%%%%%%%%%%%%%%%%%
\section{High order case\label{sec:high}}
We fix a polynomial degree $r$ with $r \geq 1$. We will define a method that properly reconstructs $k$-forms that are polynomials of degree $r$. More precisely we will define spaces $A^k(T)$ that contain $\Poly^r \Lambda^k(T)$ and such that the spaces $\Pminus^r \Lambda^{d'-k}(T')$ (with $d' = \dim T'$) constitute unisolvent degrees of freedom. The crucial point, in practice, is that we can compute the more general moments against $\Poly^r \Lambda^{d'-k} (T')$ explicitely (just removing the minus).

\begin{remark}[Link with low order case] To a large extent the low order case treated previously would correspond to $r=0$. However the theory is different on some points, because the spaces $\Pminus^0 \Lambda^k(T)$ (which look trivial but are the degrees of freedom of the Whitney forms) lack some of the properties of $\Pminus^r \Lambda^k(T)$, for $r \geq 1$.
\end{remark}
  
We set:
\begin{equation}
  U^k(T) = \{u \in \Pminus^{r+1} \Lambda^k(T) \ : \ u \perp \Poly^{r} \Lambda^k(T)\}.
\end{equation}
Orthogonality is with respect to the $\rmL^2$ product. For instance:
\begin{equation}
  U^0(T) = \{u \in \Poly^{r+1} (T) \ : \ u \perp \Poly^{r}(T)\}.
\end{equation}
In particular this space does not contain any non-zero constant function. We also notice that for $d = \dim T$:
\begin{equation}
  U^{d}(T) = 0.
\end{equation}
Then we change the previous definition we have of $W^k(T)$ to the following:
\begin{equation}\label{eq:defw}
  W^k(T) = \rmd U^{k-1}(T) + U^k(T).
\end{equation}
This entails:
\begin{align}
  W^0 & = U^0,\\
  W^d & = \rmd U^{d-1}.
\end{align}

\begin{proposition} We have:
\begin{itemize}  
\item  The decomposition (\ref{eq:defw}) is $\rmL^2$-orthogonal.

\item The following map is bijective:
  \begin{equation}
    \mapping{U^{k-1} \times U^k}{W^k(T),}{(u', u)}{\rmd u' + u.}
    \end{equation}
 
\item The sequence $W^\bs(T)$ is exact, including at indices $0$ and $d$.
\end{itemize}
\end{proposition}

\begin{proof}
\tpoint  Let $u' \in U^{k-1}(T)$ then $\rmd u' \in \Poly^{r} \Lambda^k(T)$ so $\rmd u' \perp U^k(T)$.

\tpoint  We also notice that if $\rmd u' = 0$ then $u' \in \Poly^{r} \Lambda^{k-1}(T)$ so $u' = 0$.

\tpoint  Put now $w = \rmd u' + u$, with $u' \in U^{k-1}(T)$ and $u \in U^k(T)$. If $\rmd w = 0$ then $\rmd u = 0$. Then $u \in \Poly^{r} \Lambda^k(T)$, so $u =0$.  So $w = \rmd u' \in \rmd W^{k-1}(T)$.
\end{proof}

\begin{proposition} We have direct sum decompositions:
  \begin{equation}
    \Pminus^{r+1} \Lambda^k(T) = \Pminus^r \Lambda^k(T) \oplus W^k(T).\label{eq:pmpmw}
  \end{equation}
  \end{proposition}
\begin{proof}
  \tpoint We make a preliminary observation. \\
Referring to \cite[Theorem 3.1]{ArnFalWin06},  we write, for $u \in \Poly^{r+1}\Lambda^k(T)$:
  \begin{equation}
    (\rmd \kappa_T + \kappa_T \rmd)u = (k + r+1) u + Q_r u,\textrm{with } Q_r u\in \Poly^{r}\Lambda^k(T).  
  \end{equation}
  Suppose now that $u \in U^k$ and that $\rmd u \in \Pminus^r \Lambda^{k+1}(T)$. We have:
  \begin{equation}
    (k + r + 1) u = \rmd \kappa_T u + \kappa_T \rmd u - Q_r u.
  \end{equation}
  The three terms on the right side are in $ \Poly^{r}\Lambda^k(T)$. Hence $u=0$.
  
  \tpoint  We now show that the sum is direct.\\
  Let $u' \in U^{k-1}(T)$ and $u \in U^k(T)$ and suppose $w = \rmd u' + u \in  \Pminus^r \Lambda^k(T)$. Then $\rmd u \in \Poly^{r-1}\Lambda^{k+1}(T)$ so $\rmd u = 0$, so $u= 0$. Then $\rmd u' \in  \Pminus^r \Lambda^k(T)$ so $u'=0$.
  
  \tpoint We now show that $\Pminus^{r+1} \Lambda^k(T) \subseteq \Pminus^r \Lambda^k(T) \oplus W^k(T)$.\\
  Suppose that $v \in  \Pminus^{r+1} \Lambda^k(T)$. Write first:
  \begin{equation} \label{eq:decvuv}
  v = u + v', \textrm{with } u \in U^k \textrm{ and } v' \in \Poly^{r}\Lambda^k(T).
  \end{equation}
  Next, as a variant of the above decompositon, write:
  \begin{equation}
  (\rmd \kappa_T + \kappa_T \rmd)v' = (k + r) v' + Q_{r-1} v',\textrm{with } Q_{r-1} v'\in \Poly^{r-1}\Lambda^k(T),
  \end{equation}
  and furthermore:
  \begin{equation}
  \kappa_T v' = u' + v'', \textrm{with } u' \in U^{k-1} \textrm{ and } v'' \in \Poly^{r}\Lambda^{k-1}(T).
  \end{equation}
  We now have:
  \begin{equation}
  (k+r)v' = \rmd u' + \rmd v'' + \kappa_T \rmd v' - Q_{r-1} v'.
  \end{equation}
The first term on the right side is in $\rmd U^{k-1}$ and the three others are in $\Pminus^r \Lambda^k(T)$. Inserting this in (\ref{eq:decvuv}) we are done.
  \end{proof}
  

We keep the previous definition of $Z^k_0(T)$, given the new definition of $W^{d-k}(T)$, viz.:
\begin{equation}
Z^k_0(T)  = \{u \in X^k_0(T) \ : \ \forall w \in W^{d-k}(T) \quad \int_T u \wedge w = 0\}.
\end{equation}
We remark that $Z^\bs_0(T)$ has the same exactness property as in Proposition \ref{prop:zexact}, since $W^\bs(T)$ is exact. We similarly define $Y^k(T)$ given the new definition of $Z^k_0(T)$, viz.:
\begin{equation}
  Y^k(T) = \{ u \in Z^k_0(T) \ : \ u \perp \rmd Z^{k-1}_0(T) \}.
\end{equation}

We define $G^k(T)$ as follows:
\begin{equation}
  G^k(T) = \{ (f, g) \in \Poly^{r-1} \Lambda^{k}(T) \times \Poly^{r-1} \Lambda^{k-1}(T) \ : \ \rmd^\star f = 0,\ \rmd^\star g = 0\}.
\end{equation}
A motivating observation for this definition is that for $u \in \Poly^{r} \Lambda^{k}(T)$ we have:
\begin{equation}
  (\rmd^\star \rmd u, \rmd^\star u) \in G^k(T).
\end{equation}
Concerning $k = d$ we note that $\rmd^\star f = 0$ means that $f \in \Poly^0 \Lambda^d(T) \approx \bbR$.

  
We define $B^k(T)$ to be the subspace of $X^k(T)$ consisting of fields $u \in X^k(T)$ such that the following system of PDEs holds, for some $(f,g) \in G^k(T)$:
% polynomial differential forms $f$ and $g$ of polynomial degree $r-1$, which are coclosed in the sense that $\rmd^\star f = 0$ and $\rmd^\star g = 0$:
\begin{align}\label{eq:defApde}
  \int_T \rmd u \cdot \rmd u' + \int_T \rmd u' \wedge a + \int_T u' \wedge b + \int_T u' \cdot \rmd v & = \int_T f \cdot u',\\
  \int_T \rmd u \wedge a' =0 ,\ \int_T u \wedge b' & = 0 ,\\
  \int_T u \cdot \rmd v' & = \int_T g \cdot v',
\end{align}
for all $u' \in X^k_0(T)$, for some $a \in U^{d-k-1}(T)$, some $b \in U^{d-k}(T)$, and some $v \in Y^{k-1}(T)$,  for all $v' \in Y^{k-1}(T)$, all $a' \in U^{d-k-1}(T)$ and all $b' \in U^{d-k}(T)$.

\begin{remark}[Old FES interpretation of VEM and DDR]
Removing the finite dimensional space of constraints represented by $a,a'$ and $b,b'$, and allowing general $v,v'$, we obtain the harmonic extension used in \cite{ChrGil16}. Indeed the equations considered there were:
\begin{align}
  \rmd^\star \rmd u + \rmd v & = f,\\
  \rmd^\star u & = g.
\end{align}
for some $(f,g) \in G^k(T)$. When $k = d$ the top row is replaced by $\int_T u = \int_T f$.

There, we proved that the spaces of solutions to these PDEs define a FES which is \emph{compatible} in the sense of \cite[Definition 5.12]{ChrMunOwr11} and that the spaces $\Pminus^r \Lambda^k(T)$ give unisolvent degrees of freedom. 

From this point of view, the problem was to show that the discrete $\rmL^2$ product, as defined in \cite{DiPDroRap20,BonDiPDroHu25}, is \emph{consistent} with the true $\rmL^2$ product, on this FES, in the sense of Strang's lemmas on variational crimes. We could not settle this question. This paper provides a modification of the PDEs leading to an alternative approach: the PDE is more elaborate but consistency is easy to prove, due to Proposition \ref{prop:hmom} below. 
\end{remark}

For the case $k =0$ of the definition of $B^k(T)$ we have that the system becomes:
\begin{align}
  \int_T \rmd u \cdot \rmd u' + \int_T \rmd u' \wedge a & = \int_T f \cdot u',\\
  \int_T \rmd u \wedge a' & =0 .
\end{align}
For some $a \in U^1(T)$ and all $a' \in U^1(T)$.

Concerning the case $k = d$ we have:
\begin{proposition}
  We have:
  \begin{equation}
    B^d(T) = \Poly^r\Lambda^d(T).
  \end{equation}
\end{proposition}
\begin{proof}
The system becomes:
 \begin{align}
   \int_T u' \wedge b + \int_T u' \cdot \rmd v & = \int_T f \cdot u', \label{eq:firstbd}\\
  \int_T u \wedge b' & = 0 , \label{eq:secondbd}\\
  \int_T u \cdot \rmd v' & = \int_T g \cdot v', \label{eq:lastbd}
\end{align} 
 for all $u' \in X^d(T)$, some $b \in U^0(T)$ and all $b' \in U^{0}(T)$ , some $y \in Y^{d-1}(T)$ and all $y' \in Y^{d}(T)$.



 Let $\overline u$ be the $\rmL^2$ projection of $u$ on $\Poly^r\Lambda^d(T)$. Due to the second equation (\ref{eq:secondbd}) we have:
 \begin{equation}
   \int_T (u - \overline u) \wedge b' = 0,
 \end{equation}
 for all $b' \in \Poly^{r+1}(T)$. We can choose $t \in X^{d-1}_0(T)$ such that $\rmd t = u- \overline u$ and for all $c' \in \Pminus^{r+1} \Lambda^1(T)$:
 \begin{equation}
   \int_T t \wedge c' =0.
 \end{equation}
 In particular $t \in Z^{d-1}_0(T)$.
 
 
 Since $\rmd^\star g = 0$, the last equation (\ref{eq:lastbd}) is actually true for any $v' \in Z^{d-1}_0(T)$.  Using it for $v' = t$ we get:
 \begin{equation}
   \int_T | \rmd t|^2 = \int_T g \cdot t = \pm \int_T t \wedge \hs g = 0.
 \end{equation}
 We get $u = \overline u$, so $u \in \Poly^r\Lambda^d(T)$.

 This is also sufficient for the existence of a corresponding $(f,g) \in G^d(T)$, since one can choose $f= 0$ and $g = \rmd^\star u$, and remark that $u$ solves the system with $b= 0$ and $v = 0$.
\end{proof}

 \begin{proposition}
   For $k= d$ the system (\ref{eq:defApde}) can be solved iff $f = 0$. In that case $u$ is determined up to a constant.
 \end{proposition}
 \begin{proof}
The first equation can be rewritten:
 \begin{equation}
   \pm \hs b + \rmd v = f.
 \end{equation}
 Integrating this we get $\int_T f =0$. Since $f$ is constant (by the condition $\rmd^\star f =0$), we get $f = 0$.

 On the other hand, if $f =0$, we can find $u \in \Poly^r \Lambda^d(T)$ such that $\rmd^\star u = g$. Then $u$ solves the system with $b= 0$ and $v = 0$.

 We notice that the constants solve the homogeneous problem. Suppose now that $u$ solves the homogeneous problem and $\int_T u  =0$. Since $u \in Z^d(T)$ from the second equation, we have $u \in \rmd Z^{d-1}_0(T) = \rmd Y^{d-1}(T)$. Using $v' \in Y^{d-1}(T)$ such that $u = \rmd v'$ in the last equation gives $u= 0$.
 \end{proof}
 
\begin{proposition}\label{prop:pdeext}
For $k < d$  the above PDE system is wellposed and has a unique solution $u$, given boundary data for $u$ and any right sides $f\in \rmL^2 \Lambda^k(T)$ and $g \in \rmL^2 \Lambda^{k-1}(T)$ with $\rmd^\star f = 0$ and $\rmd^\star g = 0$.
\end{proposition}
\begin{proof}
  As before (by abstract arguments involving compactness) the inf-sup condition and the coercivity on the kernel hold, for this saddlepoint problem.
\end{proof}

We define $A^k(T)$ as before from $B^k(T)$ viz.:
\begin{equation}
  A^k(T) = \{u \in X^k(T) \ : \ \forall T' \subseteq T \quad u|_{T'} \in B^k(T') \}.
\end{equation}
It is understood that $T' \subseteq T$ means that $T' \in \calT$ is a face of $T$.


We proceed to show that the spaces $A^k(T)$ constitute a FES with some very good properties.

\begin{proposition}[FES]\label{prop:highclosed}
The family of spaces $A^k(T)$ is closed under traces and the exterior derivative.
\end{proposition}
\begin{proof}
  As before, given $u \in B^k(T)$ we prove that $\rmd u \in B^{k+1}(T)$. Given $u'\in Y^k(T)$ we have, since the Lagrange multipliers cancel on $Y^k(T)$:
  \begin{align}
    \int_T \rmd u \cdot \rmd u' = \int_T f \cdot u'.
  \end{align}
  The two finite dimensional constraints required on $\rmd u$ are also satisfied.
\end{proof}

 \begin{proposition}[Approximation]
The space $ A^k(T)$ contains $\Poly^r \Lambda^k(T)$.
 \end{proposition}
 \begin{proof}
   Given $u \in \Poly^r \Lambda^k(T)$ we define $f = \rmd^\star \rmd u$ and $g = \rmd^\star u$.
   Then the PDE system is solved with zero Lagrange multipliers.

   This holds also on the faces of $T$.
 \end{proof}
 

 
 \begin{proposition}
For $k < d$ the map $G^k(T) \to A^k_0(T)$ obtained from solving (\ref{eq:defApde}) with homogeneous boundary conditions for $u$ is injective.
 \end{proposition}

 \begin{proof}
   Suppose that the solution to (\ref{eq:defApde}) with $u|_{\partial T} = 0$ is $u = 0$. We then get:
    \begin{align}
  \int_T u' \wedge w + \int_T u' \cdot \rmd v & = \int_T f \cdot u',\\
  0 & = \int_T g \cdot v',
\end{align} 
for all $u' \in X^k_0(T)$, for some $w\in W^{d-k}(T)$, all $w'\in W^{d-k}(T)$, some $v \in Y^{k-1}(T)$ and all $v' \in Y^{k-1}(T)$.

With $u' \in Y^k(T)$ we get:
\begin{equation}
  \int_T f \cdot u' =0.
\end{equation}
This also holds for $u' = \rmd v'$ for $v' \in Y^{k-1}(T)$. We deduce that for all $u' \in Z^k_0(T)$:
\begin{equation}
  \int_T f \cdot u' =0.
\end{equation}
By a density argument with respect to $\rmL^2$ norm, this can be extended to any $u' \in \rmL^2 \Lambda^k(T)$ such that:
\begin{equation}
\forall w' \in W^{d-k}(T) \quad  \int_T u' \wedge w' = 0, \textrm{ i.e. }\ u' \perp \hs W^{d-k}(T).
\end{equation}
This implies that $f \in \hs W^{d-k}(T)$.  But in view of (\ref{eq:pmpmw}) we have $\hs W^{d-k}(T) \cap \Poly^{r-1} \Lambda^k(T) = 0$, so we get $f =0$. 

The proof that $g = 0$ is identical.
 \end{proof}
 

 
\begin{proposition}[Degrees of freedom]\label{prop:hdof}
  Any element $u \in A^k(T)$ is uniquely determined by the following degrees of freedom, where $T'$ are faces of $T$:
\begin{equation}
  \int_{T'} u \wedge e, \quad e \in \Pminus^r \Lambda^{d'-k}(T'),\ d' = \dim T'.
\end{equation}

\end{proposition}

\begin{proof}
  The case $k = d$ is immediate. We proceed with the case $k <d$. Recall the notation that $A^k_0(T) = A^k(T) \cap X^k_0(T)$. We prove  that the following bilinear form on $A^k_0(T) \times \Pminus^r \Lambda^{d-k}(T)$ is invertible, in several steps:
  \begin{equation}\label{eq:dofa}
    (u, e) \mapsto \int_T u \wedge e.
  \end{equation}

  
\tpoint  For $u \in A^k_0(T)$ we obtain $u \in Z^k_0(T)$ from the finite dimensional constraints, and, for all $u' \in Z^k_0(T)$, and all $v' \in Y^{k-1}(T)$:
  \begin{align}
    \int_T \rmd u \cdot \rmd u' + \int_T u' \cdot \rmd v & = \int_T f \cdot u',\\
    \int_T  u \cdot \rmd v' & = \int_T g \cdot v'.\label{eq:ytest} 
    \end{align}
  We use $u' = \rmd v$ and use $\rmd^\star f =0$ to deduce $\rmd v = 0$ so $v= 0$.

 \tpoint  Suppose now that for all $e \in \Pminus^r \Lambda^{d-k}(T)$:
  \begin{equation}
    \int_T u \wedge e = 0.
  \end{equation}
Using (\ref{eq:pmpmw}) we then get, actually, that for all $e \in \Pminus^{r+1} \Lambda^{d-k}(T)$:
\begin{equation}\label{eq:uteste}
    \int_T u \wedge e = 0.
  \end{equation}
In particular:
\begin{equation}
  \int_T f \cdot u = 0.
\end{equation}
From the choice $u' = u$ in the first equation, this gives $\rmd u = 0$. For the case $k= 0$ we then deduce $u = 0$.


\tpoint For $0< k <d$ we proceed as follows.

If, in addition $ k> 1$, we remark that for any $v' \in Z^{k-1}_0(T)$ such that $\rmd v' = 0$ there is $v'' \in Z^{k-2}_0(T)$ such that $v' = \rmd v''$. For such $v'$ we have:
\begin{align}
  \int_T u \cdot \rmd v' & = 0,\\
  \int_T g \cdot v' & = \int_T g \cdot \rmd v'' = 0.
\end{align}
Combining this with (\ref{eq:ytest}) we conclude that for all $v' \in Z^{k-1}_0(T)$:
\begin{equation}\label{eq:ztest}
\int_T  u \cdot \rmd v'  = \int_T g \cdot v'. 
\end{equation}
If $ k =1 $ there is no $v' \in Z^{k-1}_0(T)$ such that $\rmd v' =0$, so the above conclusion also holds.

We introduce $\tilde Z^k_0(T)$ defined by:
\begin{equation}
  \tilde Z^k_0(T) = \{ u \in X^k_0(T) \ : \ \forall e \in \Pminus^{r+1} \Lambda^{d-k}(T) \quad \int_T u \wedge e = 0\}.
\end{equation}
We have $\tilde Z^k_0(T) \subseteq Z^k_0(T)$. The complex $\tilde Z^\bs_0(T)$ is exact, by the snake lemma and the exactness properties of $\Pminus^{r+1} \Lambda^{\bs}(T)$. It is a variant of Proposition \ref{prop:zexact} but now exactness of $\tilde Z^\bs_0(T)$ holds also at index $k= d$. 

Since $\rmd u = 0$ and (\ref{eq:uteste}) holds we deduce that there is $t \in \tilde Z^{k-1}_0(T)$ such that $\rmd t = u$. This entails that for all $e \in \Pminus^{r+1} \Lambda^{d-k+1}(T)$:
\begin{equation}
    \int_T t \wedge e = 0.
  \end{equation}
Then we have, setting $u = \rmd t$ and $v' = t$ in (\ref{eq:ztest}):
\begin{equation}
\int_T  \rmd t \cdot \rmd t  = \int_T g \cdot t = \int_T t \wedge \hs g = 0. 
\end{equation}
This gives $u = 0$.


We have proved that the bilinear form (\ref{eq:dofa}) defines an injective map:
\begin{equation}
  A^k_0(T) \to  \Pminus^r \Lambda^{d-k}(T)^\star.
\end{equation}

\tpoint We proceed to prove that the two spaces above have the same dimension.

For $k < d$, we have from \cite[Equation (2.28)]{ChrGil16}:
\begin{equation}
  \dim G^k(T) = \dim \Pminus^r \Lambda^{d-k}(T).
  \end{equation}
Recall also the injectivity of the map $G^k(T) \to A^k_0(T)$, proved above. We get:
\begin{equation}
\dim G^k(T) = \dim A^k_0(T) = \dim  \Pminus^r \Lambda^{d-k}(T).
\end{equation}

For $k = d$ we have more simply:
\begin{equation}
  \dim A^d(T) = \dim \Poly^r \Lambda^d(T) = \dim \Pminus^r \Lambda^0(T).
\end{equation}

% According to \cite[Equations (2.28) and (2.29)]{ChrGil16} $\dim G^k(T) = \dim \Pminus^r \Lambda^{d-k}(T)$.

It follows that the bilinear form (\ref{eq:dofa}) is invertible.

\tpoint Now  it follows from Proposition \ref{prop:pdeext} that the FES \emph{admits extensions} in the sense of \cite[Definition 2.3]{ChrRap16} (or \cite[Definition page 75]{ChrMunOwr11}). The unisolvence of the degrees of freedom follows then from a general induction argument given for FES in \cite[Proposition 2.5]{ChrRap16} (or \cite[Proposition 5.35]{ChrMunOwr11}).
\end{proof}




\begin{proposition}[Computation of moments]\label{prop:hmom}
  For elements $u \in A^k(T)$ its moments against $\Poly^r\Lambda^{d-k}(T)$ can be computed inductively from the degrees of freedom.
  \end{proposition}
\begin{proof}
\tpoint Let $u \in A^k(T)$.  Suppose that $w \in \Poly^r\Lambda^{d-k}(T)$ and that we want to compute:
  \begin{equation}
    \int_T u \wedge w.
  \end{equation}
  We write as before:
  \begin{equation}
    (d- k +r) w = \rmd \kappa_T w + \kappa_T \rmd w - Q_{r-1} w.
  \end{equation}
  On the right side the two last terms are degrees of freedom, so we concentrate on the first. We have:
  \begin{equation}\label{eq:uwk}
    \int_T u \wedge \rmd \kappa_T w = \pm \int_{\partial T} u \wedge \kappa_T w \pm \int_T \rmd u \wedge \kappa_T w.
  \end{equation}
  
  \tpoint As an intermediate step, we now treat the case where $\rmd u = 0$. Then there is only the first term on the right side. It reduces to terms of the following form, where $T'$ is a codimension one face of $T$:
  \begin{equation}\label{eq:comput}
    \int_{T'} u \wedge (\kappa_T - \kappa_{T'}) w + \int_{T'} u \wedge (\kappa_{T'} w - P_r \kappa_{T'} w) +  \int_{T'} u \wedge P_r \kappa_{T'} w.
  \end{equation}
  Here $P_r$ is $L^2$ projection onto polynomial differential forms of polynomial degree $r$. Since $P_r$ is applied to polynomial forms (or degree at most $r+1$) on a polytope, it is \emph{computable}.
  
  In (\ref{eq:comput}) the first and last terms are moments against $\Poly^r\Lambda^{d-1-k}(T')$, so amena\-ble to an induction step. The middle term is zero, since it is a moment against $U^{d-1}(T')$ representing a constraint.

  This settles the intermediate step: we can compute the moments of $u$ against polynomials of degree $r$ in the case $\rmd u = 0$, from the dofs of $u$, by induction on dimension.

  \tpoint In the general case we remark that $\rmd \rmd u = 0$ and that the degrees of freedom of $\rmd u$ are computable from those of $u$. Applying the intermediate step to $\rmd u$, its moments against polynomials of degree $r$ can be computed. Going back to (\ref{eq:uwk}) and looking at the last term, we write:
  \begin{equation}
    \int_T \rmd u \wedge \kappa_T w = \int_T \rmd u \wedge (\kappa_T w - P_r \kappa_T w) + \int_T \rmd u \wedge P_r \kappa_T w.
  \end{equation}
  On the right side the first term is zero by construction, since  $\kappa_T w - P_r \kappa_T w\in U^{d-k-1}$. The last term is computed by the above intermediate step.

  In (\ref{eq:uwk}) we are left with the first term on the right side, for which (\ref{eq:comput}) still applies. Again, this provides the beginnings of an induction procedure. 
\end{proof}


\begin{remark} We now relate some of the above results to the vocabulary of the FES framework.

  That the family of spaces $A^k(T)$ is closed under traces and the exterior derivative, as proved in Proposition \ref{prop:highclosed}, is the definition of a FES. However, not all FES produce global spaces that are suitable for Galerkin methods. Two additional conditions are important.

  The first condition, referred to as \emph{softness}, has several equivalent formulations, see especially \cite[Propositions 2.5 and 2.6]{ChrRap16}. One requires that discrete fields on the boundary of a cell, that satisfy a compatibility condition with respect to faces inside the boundary, can be extended to a discrete field on the cell. Here, this follows from Proposition \ref{prop:pdeext}. This condition turns out to be equivalent to the existence of unisolvent degrees of freedom which we have also proved, in Proposition \ref{prop:hdof}.

  The second condition, referred to as local exactness, is that for each cell $T$, the sequence $A^\bs(T)$ resolves $\bbR$. Here also, many stategies are possible. Given that softness holds, one can prove instead that the sequences $A^\bs_0(T)$ are exact, except at index $\dim T$ where the cohomology is determined by the integral, see \cite[Proposition 2.3]{ChrRap16}. This seems easier in this particular case. But even easier is to deduce the sequence exactness from a property of the degrees of freedom as in \cite[Proposition 2.7 and 2.8]{ChrRap16}. From that perspective, if $u\in A^k(T)$ and we let $v$ be a degree of freedom acting on $\rmd u \in A^{k+1}(T)$, we write:
  \begin{equation}
    \int_T \rmd u \wedge v = \int_{\partial T} u \wedge v - (-1)^k \int_T u \wedge \rmd v.
  \end{equation}
  The right side only involves degrees of freedom acting on $u$. This identity is used to assert that the discrete exterior derivative is computable in the chosen degrees of freedom and that the interpolator associated with the degrees of freedom commutes with the exterior derivative. As already indicated, a third consequence is that local exactness holds.

  The two algebraic conditions, softness and local exactness, ensure that we have a \emph{compatible} FES. When a compatible FES contains polynomials of given degree, is stable with respect to small perturbations of the mesh, and is invariant under scalings, it behaves like standard finite element spaces, such as the spaces $\Pminus^r \Lambda^k(T)$ on simplices.
\end{remark}



%%%%%%%%%%%%%%%%%%%%%%%%%%%%%%%%%%%%%%%%%%%%%%%%%%%%%%%%%%%%%%%%%%%%%%%%%%%%%%%%%%%%%%%%%%%%%%%%%%%%%%%%%%
\section{Some schemes and their analysis\label{sec:algo}}
In this section we sketch some numerical methods and elements of their analysis, relative to a choice of polynomial degree $r$. The choice $r=0$ is the lowest order case described in \S \ref{sec:low}, whereas the case $r\geq 1$ was described in \S \ref{sec:high}.

Given $u \in A^k(T)$ and based on the computation of moments against forms which are polynomials of degree $r$ obtained in Propositions \ref{prop:momconst} and \ref{prop:hmom}, one is able to inductively compute $\gamma_T u$, which is the $\rmL^2(T)$ projection of $u$ on $\Poly^r \Lambda^k (T)$.

Following VEM practice \cite{BeiBreMarRus23} one also introduces a symmetric bilinear form $s_h$ on $A^k(T)$, defined from the degrees of freedom, which is bounded and coercive in $\rmL^2$, uniformly with respect to $T$. It should scale like the $\rmL^2$ inner product (under scaling maps $x \to h_T x$, where $h_T$ is the diameter of $T$) but no consistency is required. It can be block diagonal in the degrees of freedom.

One then defines the discrete $\rmL^2$-product on $A^k(T)$ by:
\begin{equation}\label{eq:l2h}
  \langle u, u' \rangle_h = \int_T \gamma_T u \cdot \gamma_T u' + s_h( u - \gamma_T u, u' -\gamma_T u').
\end{equation}
The second term on the right side is referred to as the stabilizing term.

The crucial points are that this discrete inner product is computable, coercive and consistent. It is consistent in the sense that if at least one of $u, u'$ is in $\Poly^r \Lambda^k(T)$, it gives the exact value:
\begin{equation}\label{eq:l2exact}
\langle u , u' \rangle_h = \int_T u \cdot u'.
\end{equation}
Remark also that the exterior derivative $\rmd: A^k(T)\to A^{k+1}(T)$ is computable from the degrees of freedom on $A^k(T)$ to those on $A^{k+1}(T)$. These two properties, which are here stated in their local forms, pass to a global form. From there, numerical methods for various PDEs can be defined. Property (\ref{eq:l2exact}) make these methods amenable to an analysis based on Strang's first lemma \cite[Theorem 26.1 p. 192]{Cia91}. Such studies have been carried out in the context of VEM, as illustrated with a complete model problem in \cite[\S 3]{BeiBreMarRus23}. See also \cite{ChrHal15}.

As an example of PDE on differential forms that features many generic difficulties, we consider the Hodge-Laplace problem \cite{ArnFalWin10}. Given a $k$-form $f$ one seeks a $k$-form $u$ such that:
\begin{equation}
  - \Delta u = f, \ \textrm{with } -\Delta = \rmd^\star \rmd + \rmd\, \rmd^\star,
\end{equation}
with, say, natural boundary conditions. We suppose the domain is contractible for simplicity. One introduces the auxiliary variable $p = \rmd^\star u$ which is a $(k-1)$-form. We let $X^k$ and $X^{k-1}$ be the global spaces of $k$- and $(k-1)$-forms, which are $\rmL^2$ with exterior derivative in $\rmL^2$. The mixed formulation in terms of $(u, p) \in X^k \times X^{k-1}$ is:
\begin{equation}
\mixedright{u^\test}{X^k}{\int \rmd u \cdot \rmd u^\test + \int u^\test \cdot \rmd p}{\int f \cdot u^\test, }{p^\test}{X^{k-1}}{ \int u \cdot \rmd p^\test - \int p \cdot  p^\test}{0.}
\end{equation}
For essential rather than natural boundary conditions, one just replaces $X^k \times X^{k-1}$ by $X^k_0 \times X^{k-1}_0$.
We discretize the above problem as:
\begin{equation}
\mixedright{u^\test}{X^k_h}{\langle \rmd u , \rmd u^\test \rangle_h + \langle u^\test, \rmd p \rangle_h}{l_h(u^\test) \sim \langle f, u^\test \rangle, }{p^\test}{X^{k-1}_h}{ \langle u, \rmd p^\test \rangle_h-  \langle p ,  p^\test \rangle_h }{0.}
\end{equation}
For the convergence of such methods one establishes some stability estimates. We define:
\begin{align}
  W_h & = \{u \in X^k_h \ : \ \rmd u = 0\},\\
  V_h & = \{u \in X^k_h \ : \ \forall w \in W_h \quad \langle u, w \rangle = 0\},\\
  \widetilde V_h & = \{u \in X^k_h \ : \ \forall w \in W_h \quad \langle u, w \rangle_h = 0\}.
\end{align}

\begin{proposition}[Poincaré-Friedrichs]
  There exists $C > 0$ such that for all $u \in \widetilde V_h$ we have:
  \begin{equation}
    | u|_h \leq C | \rmd u |_h.
  \end{equation}
\end{proposition}
\begin{proof}
\tpoint  The difficulty is not so much that discrete $\rmL^2$ norms $|\cdot|_h$ appear here, than that it concerns elements of $\widetilde V_h$ rather than $V_h$. We take for granted that the result has been established for $V_h$ by FES methods involving bounded commuting projections.

\tpoint Given $u \in \widetilde V_h$, we find $v \in V_h$ such that $u$ is the projection of $v$ onto $\widetilde V_h$ with respect to $\langle \cdot , \cdot \rangle_h$. This can also be expressed by $u-v \in W_h$, i.e. $\rmd u = \rmd v$.
We then have:
\begin{align}
  |u|_h \leq |v|_h \cleq |v| \cleq |\rmd v| \cleq |\rmd v|_h = |\rmd u|_h. 
\end{align}
Here $a \cleq b$ means $a \leq C b$ for some constant $C$ independent of $h$.
\end{proof}
This, and the corresponding result in $X^{k-1}_h$ instead of $X^k_h$, is enough to establish the stability of the above discrete Hodge-Laplace problem.

\begin{remark}\label{rem:eqV}
  The situation in the above proof is that we have $u \in \widetilde V_h$ and $v \in V_h$, such that $u-v \in W_h$. The space $W_h$ is orthogonal to $V_h$ with respect to $\langle \cdot, \cdot \rangle$ and to  $\widetilde V_h$ with respect to $ \langle \cdot , \cdot \rangle_h$. Thus $v$ is the projection of $u$ onto $V_h$ with respect to $\langle  \cdot, \cdot \rangle$ and $u$ is the projection of $v$ onto $\widetilde V_h$ with espect to $\langle  \cdot, \cdot \rangle_h$. It follows that we have a uniform equivalence of norms:
  \begin{equation}
    | u | \approx | v|. 
  \end{equation}
\end{remark}

  
  

For corresponding eigenvalue problems such as those of Maxwell \cite{Bof10,ChrWin13IMA} one needs a stronger result called discrete compactness. Some further notations:
\begin{align}
  W & = \{u \in X^k \ : \ \rmd u = 0\},\\
  V & = \{u \in X^k \ : \ \forall w \in W \quad \langle u, w \rangle = 0\}.
 % \widetilde V_h & = \{u \in X_h \ : \ \forall w \in W_h \quad \langle u, w \rangle_h = 0\}.
\end{align}
Let $P_V$ be the $\rmL^2$ projection onto $V$. We suppose that we have $\rmL^2$ bounded commuting projections $\Pi_h$ onto $X^k_h$. We start with a $v \in V_h$. We notice that $v$, $P_V v$ and $\Pi_h P_V v$ have the same exterior derivative. So $v - \Pi_h P_V v \in W_h$ is $\rmL^2$ orthogonal to both $v$ and $P_V v$. We deduce: 
\begin{equation}
  |v - P_V v|^2 + |v - \Pi_h P_V v |^2 = |P_V v - \Pi_h P_V v|^2.  
\end{equation}
Hence
\begin{equation}
  |v - P_V v| \leq  |P_V v - \Pi_h P_V v|.
\end{equation}
The space $V$, which is equipped with the norm of $X^k$, encodes also vanishing coderivative. An alternative norm on $V$ is $|\rmd \cdot |$. The injection of $V$ into $\rmL^2$ is compact. One deduces that for some sequence $\epsilon_h \to 0$:
\begin{equation}
  |v - P_V v| \leq \epsilon_h | \rmd v|.
\end{equation}
This gives the vanishing gap condition (concerning the one-sided gap from $V_h$ to $V$ in $X^k$ norm), which is equivalent to discrete compactness \cite[Proposition 3.9]{ChrWin13IMA}. If $(v_h)$ is a sequence in $\widetilde V_h$ with $\rmL^2$ bounded $(\rmd v_h)$, one can get a subsequence converging in the $\rmL^2$ norm, simply by considering $(P_V v_h)$.


We now prove the vanishing gap condition also for $\widetilde V_h$.
 \begin{proposition}[Vanishing gap]
  There is a sequence $\epsilon_h \to 0$ such that when $u \in \widetilde V_h$ and $v \in V_h$ is defined (as before) by $\rmd v = \rmd u$ we have:
  \begin{equation}
    \| u - v\| \leq \epsilon_h \| u\|.
  \end{equation}
\end{proposition}
\begin{proof}
  Let $Q_h$ be the map $X^k \to X^k_h$ defined so that, for $t' \in X^k_h$:
  \begin{equation}
    \langle Q_h t,t' \rangle_h = \langle t, t'\rangle. 
  \end{equation}
  It is not a projection, but it maps $V$ to $\widetilde V_h$ and satisfies approximation properties (by stability and consistence).
  
We have:
\begin{align}
  |u - v|^2_h & = \langle u - v, Q_h P_V u - v \rangle_h \leq | u-v |_h |Q_h P_V u - v|_h.
\end{align}
We get, remarking that $P_V u = P_V v$ (they have the same exterior derivative):
\begin{align}
  |u -v| & \cleq |Q_h P_V u - v| \leq |Q_h P_V u - P_V u| + |P_V v - v|,\\
  & \cleq \epsilon_h | \rmd u|,
\end{align}
using the already established results.
\end{proof}

\begin{remark}
  This proposition can be interpreted as saying that the symmetrized gap $\delta(\widetilde V_h, V_h)$, computed in the norm $\| \cdot \|$ of $X^k$, converges to $0$ as $h \to 0$. On the other hand Remark \ref{rem:eqV} can be interpreted as saying that the symmetrized gap $\delta(\widetilde V_h, V_h)$ computed in the $\rmL^2$ norm $| \cdot|$ stays away from 1.
\end{remark}

  
\bibliography{../Bibliography/alexandria,../Bibliography/newalexandria,../Bibliography/mybibliography}{}
\bibliographystyle{plain}

\end{document}

%%%%%%%%%%%%%%%%%%%%%%%%%%%%%%%%%%%%%%%%%%%%%%%%%%%%%%%%%%%%%%%%%%%%%%%%%%%%%%%%%%%%%%%%%%%%%%%%%%%%%%%%%%%%%%%%%%%%%































